\article{若干高维模簇的 Hasse--Weil $\zeta$ 函数}{The Hasse-Weil zeta-function of some moduli varieties of dimension greater than one}{W. Casselman}{2}{141--163}

\section*{引言（Introduction）}
设 $k$ 为数域（number field），$G$ 为 $k$ 上的连通约化群（connected reductive group）。通过限制标量（restriction of scalars），不妨假设 $k=\Q$。令 $K_\R\subset G(\R)$ 为 $G(\R)$ 的中心与一个极大紧子群的乘积。则在某些条件下，可赋予 $X=G(\R)/K_\R$ 一个 $G(\R)$-不变的复结构（complex structure）；如果 $\Gamma$ 是 $G$ 的一个足够小、从而不含挠元的算术子群（arithmetic subgroup），则 $\Gamma\backslash X$ 为非奇异代数簇（nonsingular algebraic variety）。在许多情形中，可以为它构造定义在某个明确数域上的模型。加上进一步的假设，还可选取定义在特定数域 $E$ 的阿贝尔扩张上的典范模型（canonical model）；Deligne 与 Langlands 的报告或将讨论这一点。粗略地说，$E$ 由 $G$ 以及 $X$ 上的复结构确定。\footnote{原文的 AMS（MOS）主题分类（1970）：主要分类 10D20；次要分类 14G25、22E55。原文版权：1979，American Mathematical Society。}

一个自然的猜想是，典范模型的 Hasse--Weil $\zeta$ 函数是 Langlands 赋予 $G(\A)$ 的自守表示（automorphic representation）的那类 $L$ 函数的乘积。事实证明这并不正确（见 [22] 的引言），但这一想法仍很有启发。这一方向的首个结果由 Eichler 得到：当 $G=\GL_2(\Q)$、$\Gamma=\Gamma_0(N)$ 时，$\Gamma\backslash X$ 有一个定义在 $\Q$ 上的模型，并且除 Euler 乘积中的有限多个因子外，其 Hasse--Weil $\zeta$ 函数是 Hecke 在这一情形定义的 $L$ 函数的乘积。后人尤其是 Shimura、Kuga、Ihara、Deligne 和 Langlands，将此结果推广到：(a) 该 $G$ 中的其他 $\Gamma$；(b) 其他 $G$；(c) 与非平凡局部常值层（locally constant sheaf）相联系的 $\zeta$ 函数；以及 (d) 坏约化素数处的 $\zeta$ 因子。除一个例外，即 Shimura 的一些未发表工作外，这些工作都讨论维数为一的 $X$（有关例子见 [11]、[28]、[14]、[6] 和 [19]）。

应用 Deligne 关于有限域上代数簇的 $\zeta$ 函数之根的结果，可由上述结果得到 Ramanujan 猜想的一个推广。另一个推论是所讨论的 Hasse--Weil $\zeta$ 函数的函数方程（functional equation）与解析延拓（analytic continuation）；除 Shimura 的例子外，它都是以特别简单的方式与 $\GL_2$ 或四元数代数（quaternion algebra）相联系的 $\zeta$ 函数。

过去几年，Langlands 着手研究维数大于一的簇的情形；Milne 与我在讲座中将讨论他处理的最简单情形。

更确切地，令 $F$ 为 $n$ 次全实域（totally real field），$\mathfrak o_F$ 为 $F$ 的整数环，$B$ 为 $F$ 上的四元数代数，$\mathfrak o_B$ 为 $B$ 中的极大序（maximal order）。这里，四元数代数是指 $F$ 上的四维代数 $B$，且 $B\otimes\overline F=M_2(\overline F)$，其中 $\overline F$ 为 $F$ 的代数闭包。由于在这里出现的每个完备化 $F_v$ 上都存在唯一的四元数除代数，如果 $v$ 是 $F$ 的任意赋值，则 $B_v=B\otimes_FF_v$ 根据是否有零因子，或者同构于 $M_2(F_v)$，或者同构于该唯一的除代数。前一种情形称 $B$ 在 $v$ 处分裂（split），后一种称其分歧（ramified）。实际上，除有限多个赋值外，$B$ 在所有赋值处分裂；二次互反律说明这个有限数目为偶数。另一个经典结果表明，对每个由偶数个素位组成的有限集合，都存在唯一的四元数代数，恰好在那些赋值处分歧。因此，$B$ 本身为除代数，当且仅当它至少在某处分歧。若 $v$ 非 Archimedes，则 $\mathfrak o_B$ 在 $B_v$ 中的闭包为 $B_v$ 的极大紧子环；若 $B$ 在 $v$ 处分歧，则此子环唯一，否则仅在 $B_v^\times$-共轭意义下唯一。

令 $G$ 为由 $\mathfrak o_B$ 的乘法群定义的 $\Z$ 上的代数群。因此对任意环 $R$，有 $G(R)=(\mathfrak o_B\otimes R)^\times$。特别地，$G(\Q)\simeq B^\times$（典范地），而
\[
 G(\R)\simeq\GL_2(\R)^I\times(\mathbb H^\times)^J
\]
（非典范地），其中 $I$ 为 $F$ 的实赋值中 $B$ 分裂者的集合，$J$ 为其分歧者的集合。对每个 $B$ 不分歧于其上的有理有限素数 $p$，有
\[
 G(\Z_p)\simeq\prod_{\mathfrak p\mid p}\GL_2(\mathfrak o_{F,\mathfrak p}),
 \qquad G(\Q_p)\simeq\prod_{\mathfrak p\mid p}\GL_2(F_{\mathfrak p}).
\]

最简单的情形是 $B=M_2(F)$；本文将讨论除代数的情形，并作如下假设：\emph{今后假定 $B$ 为除代数。此外，假定 $B$ 在 $F$ 的所有实素位均不定，即 $J=\varnothing$。}Langlands 本人没有作这些假设，但他承认，除非这些假设成立，否则其论证尚有缺口。

令 $\Delta$ 为有限有理素数集合，包含 $F$ 或 $B$ 在其上分歧的素数；令 $K_f$ 为 $G(\Z_f)$ 中形如
\[
 K_\Delta\cdot\prod_{p\notin\Delta}G(\Z_p)
\]
的紧开子群，其中 $K_\Delta$ 为 $\prod_{p\in\Delta}G(\Z_p)$ 的紧开子群。令 $Z$ 为 $G$ 的中心，$Z_K=(Z(\A_f)\cap K_f)\cdot Z(\R)$。将 $\C^\times$ 嵌入 $\GL_2(\R)$：
\[
 a+b\sqrt{-1}\longmapsto\begin{pmatrix}a&-b\\b&a\end{pmatrix},
\]
并令 $K_\R$ 为 $(\C^\times)^I$ 在 $G(\R)$ 中的像（这是上文使用的 $K_\R$ 的连通分量）。于是 $G(\R)/K_\R$ 是 $n$ 个 $\C\setminus\R$ 的乘积。

令 $K=K_\R K_f$。集合 ${}_KS(\C)=G(\Q)\backslash G(\A)/K$ 是有限个 $n$ 维紧复解析空间的并，后文将说明这一点。如果 $K_f$ 足够小，它便是非奇异的；下文均作这一假设。

Milne 在本研究班的讲座中将证明，此空间是某个模概形（moduli scheme）${}_KS$ 的复值点集合，且 ${}_KS$ 定义在 $\Spec\Z[1/d]$ 上，在其上光滑且固有（proper），其中 $d$ 为 $\Delta$ 中素数的乘积。他还将讨论它在有限域上的点的结构。本文将说明如何用 Langlands 的 $L$ 函数表达其 Hasse--Weil $\zeta$ 函数。基本思路是：用 Selberg 迹公式（Selberg trace formula）计算 Langlands 的 $L$ 函数的 $p$ 因子；用 Milne 的结果和迹公式的一种非传统应用计算 Hasse--Weil 的 $p$ 因子，二者均对 $p\notin\Delta$ 进行。具体论证较为曲折。先说明一点：虽然当 ${}_KS$ 为一维时，可用一个同余关系（congruence relation）得到最终结果，但一般情形仅此并不足够。随后可应用 Deligne 关于 Frobenius 根的结果，得到 Ramanujan--Petersson 的一个推广；不过，关于 Hasse--Weil 函数方程的结果尚不完整。

Milne 与我的工作本质上都是在报告 Langlands 的工作。我们讨论的结果，只是他一个涉及 $B^\times$ 的子群的更复杂结果的特例。处理这个特例，避开了他称为“$L$ 不可区分性”（$L$-indistinguishability）的问题，也许他本人会谈到这个问题。他的结果在其他方面也更一般；最后一节补充讨论若干相关问题，也将介绍这些推广。

主要参考文献是概要 [21]；完整陈述与证明见 [22]。Langlands 在 Hilbert 问题讨论会上的报告 [20]，以及他在本研究班的讲座，也与此有关。

\origsection{1}{${}_KS(\C)$ 的上同调与表示（Cohomology and representations）}
\origsubsection{1.1}{}
令 $\nu:G\to\mathbb G_m$ 为约化范数（reduced norm）。\footnote{\textbf{译者注（本次校订）：}由于这里的 $G$ 已通过限制标量视为 $\Q$ 上的群，值域为 $\mathbb G_m$ 的 $\nu$ 应理解为复合 $N_{F/\Q}\circ\operatorname{Nrd}_{B/F}$。单独的约化范数取值于 $\operatorname{Res}_{F/\Q}\mathbb G_m$。}由于 $B$ 为除代数，陪集空间 $G(\Q)\backslash G^1(\A)$ 紧，其中 $G^1(\A)$ 为模同态（modulus homomorphism）$|\nu|:x\mapsto|\nu(x)|$ 的核。由于 $G(\R)$ 的连通分量 $G^{\mathrm{conn}}(\R)$ 在 $|\nu|$ 下的像为整个 $\R^{\mathrm{pos}}$，集合 $G(\Q)\backslash G(\A)/G^{\mathrm{conn}}(\R)$ 也紧。$K_f$ 在 $G(\A_f)$ 中开，因此 $G(\Q)\backslash G(\A)/G^{\mathrm{conn}}(\R)K_f$ 有限。故存在 $G(\A)$ 中元素的有限集合 $\mathcal X$，使
\[
 G(\A)=\bigcup_{x\in\mathcal X}G(\Q)xG^{\mathrm{conn}}(\R)K_f.
\]
强逼近（strong approximation）可对 $\mathcal X$ 作更具体的参数化，但这里不需要；见 [7]。

\begin{originaltheorem}{引理 1.1.1（Lemma 1.1.1）}
对任意 $x\in G(\A)$，空间
\[
 G(\Q)\backslash G(\Q)xG^{\mathrm{conn}}(\R)K_f/K_f
\]
作为 $G^{\mathrm{conn}}(\R)$-空间，同构于 $\Gamma_x\backslash G^{\mathrm{conn}}(\R)$，其中 $\Gamma_x$ 为 $G(\Q)\cap G^{\mathrm{conn}}(\R)\cdot xK_fx^{-1}$ 在 $G^{\mathrm{conn}}(\R)$ 中的像。
\end{originaltheorem}

这是因为 $G^{\mathrm{conn}}(\R)$ 显然在该空间上传递作用，而 $\Gamma_x$ 是陪集 $G(\Q)xK_f$ 的稳定子群（isotropy subgroup）。

令 $\mathcal H$ 为 $\C$ 中的上半平面，由上面得到：

\begin{originaltheorem}{命题 1.1.2（Proposition 1.1.2）}
$G^{\mathrm{conn}}(\R)$-空间 $G(\Q)\backslash G(\A)/K_f$ 同构于空间 $\Gamma_x\backslash G^{\mathrm{conn}}(\R)$（$x\in\mathcal X$）的不交并。簇 ${}_KS(\C)$ 为 $\Gamma_x\backslash\mathcal H^n$ 的不交并。
\end{originaltheorem}

可采用与 [19] 引理 2.1 的证明相同的论证，表明只要 $K_f$ 足够小，每个 $\Gamma_x$ 的作用都是自由的，且 ${}_KS(\C)$ 非奇异。\footnote{\textbf{译者注（本次校订）：}这里的自由作用应指模去中心核后的有效作用。$\Gamma_x$ 中的中心元素在 $\mathcal H^n$ 上恒等作用；当 $[F:\Q]>1$ 时，缩小 $K_f$ 仍会留下无穷多个中心单位。相应的有效算术群为下节使用的 $\widetilde\Gamma_x$。}\emph{今后均作此假设。}

\origsubsection{1.2}{}
令 $\mathcal A_0$ 为 $G(\Q)Z(\R)\backslash G(\A)$ 上的自守形式空间。它为 $G(\A)$ 的不可约、容许、酉表示的直和 $\bigoplus\pi$；这里采用通常的简略说法：在所取的有限向量空间上，实位的作用通过 Lie 代数 $\mathfrak g$ 来描述。
若 $\widetilde{\mathfrak g}$ 为 $\widetilde G(\R)=G(\R)/Z(\R)$ 的 Lie 代数，则 $\mathcal A_0$ 上的表示实际上通过 $\widetilde{\mathfrak g}$ 分解。令 $\widetilde{\mathfrak k}$ 为 $\widetilde K_\R=K_\R/(K_\R\cap Z_\R)$ 的 Lie 代数。

\begin{originaltheorem}{命题 1.2.1（Proposition 1.2.1）}
de Rham 上同调（de Rham cohomology）$H^*({}_KS(\C),\C)$ 自然同构于
\[
 \bigoplus_\pi H^*(\widetilde{\mathfrak g},\widetilde{\mathfrak k},\pi_\infty)
 \otimes\pi_f^{K_f}.
\]
\end{originaltheorem}

这里是相对 Lie 代数上同调（relative Lie algebra cohomology）。求和遍及 $\mathcal A_0$ 的所有组成表示 $\pi$，必要时按重数计。每个 $\pi$ 分解为 $\pi=\pi_\infty\otimes\pi_f$，其中 $\pi_\infty$ 为 $G(\R)$ 的不可约容许表示，$\pi_f$ 为 $G(\A_f)$ 的不可约容许表示；实位的表示仍按上述约定理解。

\begin{proof}[证明梗概]
令 $\mathcal A_0^{K_f}$ 为 $\mathcal A_0$ 中被 $K_f$ 的元素固定的函数子空间。由于 $\widetilde{\mathfrak g}$ 与 $K_f$ 交换，作为 $\widetilde{\mathfrak g}$-表示，它自然同构于 $\bigoplus\pi_\infty\otimes\pi_f^{K_f}$。因而命题等价于证明 $H^*({}_KS(\C),\C)$ 同构于 $H^*(\widetilde{\mathfrak g},\widetilde{\mathfrak k},\mathcal A_0^{K_f})$。由 1.1.2，$\mathcal A_0^{K_f}$ 是子空间 $\mathcal A_{0,x}^{K_f}$（$x\in\mathcal X$）的直和，其中每个子空间为
\[
 \mathcal A(\widetilde\Gamma_x\backslash\widetilde G^{\mathrm{conn}}(\R)),
\]
即 $\widetilde\Gamma_x\backslash\widetilde G^{\mathrm{conn}}(\R)$ 上的 $K_\R$-有限、$Z(\widetilde{\mathfrak g})$-有限函数空间；$\widetilde\Gamma_x$ 为 $\Gamma_x$ 在 $\widetilde G$ 中的像。因此命题 1.2.1 的证明归结为下述引理。
\end{proof}

\begin{originaltheorem}{引理 1.2.2（Lemma 1.2.2）}
存在自然同构
\[
 H^*(\widetilde\Gamma_x\backslash\widetilde G/\widetilde K,\C)
 \simeq H^*(\widetilde{\mathfrak g},\widetilde{\mathfrak k},
 \mathcal A(\widetilde\Gamma_x\backslash\widetilde G)).
\]
\end{originaltheorem}

实际上，对任意半单 Lie 群 $\widetilde G$、余紧的 $\widetilde\Gamma\subset\widetilde G$ 及极大紧子群 $\widetilde K$，此类结论都成立。说明如下：$X=\widetilde\Gamma_x\backslash\widetilde G/\widetilde K$ 的上同调为 de Rham 复形（de Rham complex）的上同调，其第 $m$ 项为 $X$ 上的 $C^\infty$ 的 $m$ 次微分形式空间。投影 $\widetilde\Gamma_x\backslash\widetilde G\to X$ 为主丛（principal bundle），而 $m$ 次微分形式丛由此主丛及 $\widetilde K$-空间 $\Lambda^m(\widetilde{\mathfrak g}/\widetilde{\mathfrak k})^\wedge$ 关联而得。因此，$C^\infty$ 的 $m$ 次形式对应于从 $\widetilde\Gamma_x\backslash\widetilde G$ 到 $\Lambda^m(\widetilde{\mathfrak g}/\widetilde{\mathfrak k})^\wedge$ 的某些 $C^\infty$ 函数；利用对偶关系，也可将它们视为从 $\Lambda^m(\widetilde{\mathfrak g}/\widetilde{\mathfrak k})$ 到 $C^\infty(\widetilde\Gamma_x\backslash\widetilde G)$ 的 $\widetilde K$-线性映射。换言之，$X$ 上的 de Rham 复形可以等同于第 $m$ 项为
\[
 \Hom_{\widetilde{\mathfrak k}}
 \bigl(\Lambda^m(\widetilde{\mathfrak g}/\widetilde{\mathfrak k}),
 C^\infty(\widetilde\Gamma_x\backslash\widetilde G)\bigr)
\]
的复形。这正是计算 $H^*(\widetilde{\mathfrak g},\widetilde{\mathfrak k},C^\infty(\widetilde\Gamma_x\backslash\widetilde G))$ 的复形的第 $m$ 项。显式计算表明，其微分也相同。因此，若把引理 1.2.2 中的 $\mathcal A(\widetilde\Gamma_x\backslash\widetilde G)$ 换成 $C^\infty(\widetilde\Gamma_x\backslash\widetilde G)$，结论成立。最后使用 Hodge 理论（Hodge theory）：每个上同调类都有唯一的调和代表，而这些调和代表在上述对应下提升为 $\mathcal A$ 中的元素。关于这一点和其他细节，见 [3] 第 IV 章。

\origsubsection{1.3}{}
虽然 1.2.1 的和中出现无限多个表示，但除有限多项外，其余项均为零。为明确哪些项有贡献，需进一步讨论容许表示的相对 Lie 代数上同调。首先，由于 $\widetilde G(\R)\simeq\operatorname{PGL}_2(\R)^I$（记号同引言），每个 $\pi_\infty$ 都分解为 $\bigotimes\pi_{\infty,\iota}$，其中 $\pi_{\infty,\iota}$ 为 $\operatorname{PGL}_2(\R)$ 的容许表示。由 Künneth 公式（Künneth formula）
\[
 H^*(\widetilde{\mathfrak g},\widetilde{\mathfrak k},
 \bigotimes_\iota\pi_{\infty,\iota})
 \simeq\bigotimes_\iota H^*(\mathfrak{sl}_2,\mathfrak{so}(2),\pi_{\infty,\iota}),
\]
因此 $G(\R)$ 的问题归结为 $\operatorname{PGL}_2(\R)$ 的问题。注意 $\operatorname{PGL}_2$ 与 $\SL_2$ 的 Lie 代数相同。令 $C$ 为 $U(\mathfrak{sl}_2)$ 中的 Casimir 元素（Casimir element）；它属于 $U(\mathfrak{sl}_2)$ 的中心，并被极大紧子群 $O(2)$ 中心化，因而在 $\operatorname{PGL}_2(\R)$ 的任意不可约容许表示上按标量作用。

\begin{originaltheorem}{引理 1.3.1（Lemma 1.3.1）}
设 $(\pi,V)$ 为 $\operatorname{PGL}_2(\R)$ 的不可约、容许、酉表示。则
\[
 H^*(\mathfrak{sl}_2,\mathfrak{so}_2,V)\simeq
 \begin{cases}
 0,&\pi(C)\ne0,\\
 \Hom_{\mathfrak{so}_2}(\Lambda^*(\mathfrak{sl}_2/\mathfrak{so}_2),V),&\pi(C)=0.
 \end{cases}
\]
\end{originaltheorem}

证明思想是：若 $\pi$ 酉，则可在复形 $\Hom_{\mathfrak{so}_2}(\Lambda^*(\mathfrak{sl}_2/\mathfrak{so}_2),V)$ 上赋予自然内积，使 $\pi(C)=dd^*+d^*d$，其中 $d^*$ 为 $d$ 的伴随（adjoint）。这一对所有半单群都成立的观察出自 Kuga。引理立即由此推出。

对于 $\operatorname{PGL}_2(\R)$，只有三个不可约容许表示满足 $\pi(C)=0$：(a) 平凡表示 $\C$；(b) 特征 $\operatorname{sgn}(\det(g))$；(c) 一个离散系列表示（discrete series representation）$\pi_0$，可与 $\mathbf P^1(\R)$ 上的 $O(2)$-有限函数模去常值函数所得的商空间等同。结合引理 1.3.1，并计算这些表示在 $O(2)$ 上的限制，得到：

\begin{originaltheorem}{推论 1.3.2（Corollary 1.3.2）}
若 $\pi$ 为 $\operatorname{PGL}_2(\R)$ 的不可约容许酉表示，且不为上述三个之一，则 $H^*(\pi)=H^*(\mathfrak{sl}_2,\mathfrak{so}_2,\pi)=0$。对上述表示：

(a) 当 $\pi=\C$ 或 $\operatorname{sgn}(\det)$ 时，
\[
 H^m(\pi)\simeq\begin{cases}\C,&m=0,\\0,&m=1,\\\C,&m=2;\end{cases}
\]
(b) 当 $\pi=\pi_0$ 时，
\[
 H^m(\pi)\simeq\begin{cases}0,&m=0,\\\C\oplus\C,&m=1,\\0,&m=2.\end{cases}
\]
\end{originaltheorem}

实际上不需要酉性假设（见 [3]）。因此，若原来 $G(\R)$ 的不可约容许表示 $\pi_\infty$ 的上同调非平凡，它必为 $\bigotimes\pi_{\infty,\iota}$ 的形式，其中每个 $\pi_{\infty,\iota}$ 是上述三个表示之一，其上同调可据此计算。对 $\operatorname{PGL}_2(\R)$ 的不可约容许表示 $\pi$，令
\[
 m(\pi)=\begin{cases}
 0,&\pi\text{ 的上同调平凡},\\
 1,&\pi=\C\text{ 或 }\operatorname{sgn}(\det),\\
 -1,&\pi=\pi_0.
 \end{cases}
\]
若 $\pi_\infty=\bigotimes\pi_{\infty,\iota}$ 是 $G(\R)$ 的不可约容许表示，则定义 $m(\pi_\infty)=\prod m(\pi_{\infty,\iota})$。若 $\pi=\pi_\infty\otimes\pi_f$ 为 $G(\A)$ 的不可约容许表示，且 $\pi_\infty$ 在 $Z(\R)$ 上平凡，则定义
\[
 m(\pi,K)=m(\pi_\infty)\dim\pi_f^{K_f}.
\]
因此 $m(\pi,K)$ 的绝对值恰为 $\dim\pi_f^{K_f}$\footnote{\textbf{译者注（本次校订）：}此句适用于 $m(\pi_\infty)\ne0$ 的表示，即实位相对 Lie 代数上同调非零的情形；若 $m(\pi_\infty)=0$，则 $m(\pi,K)=0$，而 $\pi_f^{K_f}$ 仍可能非零。}，其符号反映该表示对上同调贡献的奇偶性。

最后指出，应用强逼近定理可证明：若 $\pi=\bigotimes_v\pi_v$（$v$ 遍及 $F$ 的所有赋值）是 $G(\A)$ 的不可约容许自守表示，且 $\pi_v$ 在某个 $B$ 分裂的素位处是一维的，则它在所有素位处都是一维的。特别地，若 $\pi_\infty$ 的一个因子是一维的，则所有因子都是一维的。因此重新得到 Matsushima--Shimura [23] 的结果：${}_KS(\C)$ 中超出一维表示贡献的上同调集中在中间次数。

\origsection{2}{主要定理及若干推论（The main theorem and some consequences）}
\origsubsection{2.1}{}
先回顾 $G$ 的 $L$ 群。首先，若以最直接的方式将 $G$ 视为 $F$ 上的群，即对 $F$ 的任意扩域 $F'$，有 $G(F')\simeq(B\otimes_FF')^\times$，则其 $L$ 群就是 ${}^LG^0=\GL_2(\C)$ 与 $\Gal(\overline F/F)$ 的直积，因为它是 $\GL_2(F)$ 的一个内扭曲（inner twisting）。由于 $\Q$ 上的群 $G$ 由此群限制标量而得，其 $L$ 群在某种意义上也由前者诱导而来：它是 $\Gal(\overline\Q/\Q)$ 与 ${}^LG^0=\GL_2(\C)^I$ 的半直积，其中 $I$ 为 $F$ 到 $\overline\Q$ 的嵌入集合，Galois 群通过置换因子作用。由于 $\overline\Q$ 可视为 $\C$ 的子域，这个 $I$ 本质上与前文相同。在本文的讨论中，可以用 $\Gal(F_{\mathrm{norm}}/F)$ 替换 $\Gal(\overline\Q/\Q)$，其中 $F_{\mathrm{norm}}$ 为包含 $F$ 且在 $\Q$ 上正规的最小扩域。还可注意到，只要 $F$ 在某个有理素数处非分歧，该正规闭包也在该处非分歧。若有理素数 $p$ 在 $F$ 中非分歧，$\Phi$ 为 $\Gal(F_{\mathrm{norm}}/F)$ 中位于 $p$ 之上的 Frobenius，则局部 $L$ 群 ${}^LG_{\Q_p}={}^LG_p$ 可等同于 ${}^LG_\Q$ 中在 $\Gal(F_{\mathrm{norm}}/F)$ 下的像属于 $\Phi$ 生成的循环子群的子群。因此有正合列
\[
 1\longrightarrow\GL_2(\C)^I\longrightarrow{}^LG_p
 \longrightarrow\langle\Phi\rangle\longrightarrow1.
\]
\noindent\emph{译注：所供扫描本本段数处写作 $\Gal(F_{\mathrm{norm}}/F)$，而定义限制标量的 $L$ 群时所需的有限 Galois 群通常为 $\Gal(F_{\mathrm{norm}}/\Q)$；本段保留扫描符号，后一段亦印出以 $\Q$ 为基域的形式。}

关于 $L$ 群的全部内容，参见 [2]。

为定义与自守表示相联系的 $L$ 函数，还须引入 ${}^LG$ 的有限维表示。对每个 $G$，我们只关心一个这样的表示 $\rho$，定义如下：其空间为若干个 $\C^2$ 的张量积，对 $I$ 的每个元素取一个因子。群 ${}^LG^0=\GL_2(\C)^I$ 通过各因子上的标准表示作用，$\Gal(\overline\Q/\Q)$（或 $\Gal(F_{\mathrm{norm}}/\Q)$，没有区别）则通过置换因子作用。由于 ${}^LG$ 为这两个群的半直积，这确实定义一个表示。其维数为 $2^n$。例如当 $F=\Q$ 时，${}^LG=\GL_2(\C)$，$\rho$ 就是标准表示本身。

这一 $\rho$ 的选择来自 Langlands 对 Shimura 簇（Shimura variety）的总体构想；参见 [22] 引言对此的讨论。

\origsubsection{2.2}{}
设 $X$ 是 $\Z[1/d]$ 上的任意光滑固有概形（smooth proper scheme），其中 $d\ge1$，$p$ 为不整除 $d$ 的素数。则 $X$ 在 $\F_p$ 上的 $\zeta$ 函数定义为 $p^{-s}$ 的有理函数 $Z_p(s,X)$，使至少形式地有
\[
 \log Z_p(s,X)=\sum_{m=1}^\infty\frac1{mp^{ms}}\#X(\F_{p^m}).
\]
由平展理论（étale theory），这与利用 $l$ 进上同调（$l$-adic cohomology）的定义一致：
\[
 Z_p(s,X)=\prod_{i=0}^{2\dim X}
 \det(I-\Phi p^{-s})_{H^i(X,\Q_l)},
\]
\noindent\emph{译注：所供扫描本的上述乘积未见交错指数。通常的上同调表达式中，第 $i$ 个行列式因子的指数为 $(-1)^{i+1}$；这里保留扫描本所示形式。}

其中 $\Phi$ 为几何 Frobenius（geometric Frobenius）。至少对不整除 $d$ 的素数处的因子而言，Hasse--Weil $\zeta$ 函数 $Z(s,X)$ 为这些 $Z_p(s,X)$ 的乘积。

Langlands 的主要定理是：

\begin{originaltheorem}{定理 2.2.1（Theorem 2.2.1）}
除 $\Delta$ 中素数对应的因子外，${}_KS$ 的 Hasse--Weil $\zeta$ 函数等于
\[
 \prod_\pi L(s-n/2,\pi,\rho)^{m(\pi,K)},
\]
其中乘积遍及作为 $\mathcal A_0$ 的组成表示出现的所有 $\pi$。
\end{originaltheorem}

乘积中的 $\pi$ 应按其重数计入。不过，[16] 第 16 节将 $G$ 上自守形式与 $\GL_2$ 上自守形式联系起来的定理，联合该文第 9 节关于 $\GL_2$ 的结果，实际上说明每个 $\pi$ 恰好出现一次。

此乘积中，除有限多个 $\pi$ 外，均有 $m(\pi,K)=0$，这是因为 ${}_KS(\C)$ 的上同调是有限维的。

实际要证明的是纯局部结果：对 $p\notin\Delta$，$Z(s,{}_KS)$ 与上述 $L$ 函数乘积的 $p$ 因子相等。对于每个满足 $m(\pi,K)\ne0$ 的 $\pi=\bigotimes_v\pi_v$，其 $p\notin\Delta$ 处的分量 $\pi_p$ 非分歧，故对应于局部 $L$ 群 ${}^LG_p$ 中的元素 $g(\pi_p)$。$L(s-n/2,\pi,\rho)$ 的 $p$ 因子为
\[
 \det\bigl(I-\rho(g(\pi_p))/p^{s-n/2}\bigr)^{-1}.
\]
展开后，定理化为形式等式
\[
 \begin{aligned}
 \sum_{m=1}^\infty\frac{\#{}_KS(\F_{p^m})}{mp^{ms}}
 =\sum_{\pi\text{ 在 }\mathcal A_0\text{ 中}}
 m(\pi,K)\sum_{m=1}^\infty
 \frac{p^{mn/2}}{mp^{ms}}\operatorname{trace}\rho(g(\pi_p)^m),
 \end{aligned}
\]
即当 $p\notin\Delta$、$m\ge1$ 时，系数满足
\begin{equation}
 \#{}_KS(\F_{p^m})=
 \sum_{\pi\text{ 在 }\mathcal A_0\text{ 中}}
 m(\pi,K)p^{mn/2}\operatorname{trace}\rho(g(\pi_p)^m).
 \tag{2.1}
\end{equation}
下面将证明定理的这一等价形式。本讲座仅在最简单的情形，即 $p$ 在 $F$ 中分裂时，给出完整证明。

\origsubsection{2.3}{}
开始证明之前，先给出一个例子与若干推论。$n=2$ 是这里首次出现的新情形。如前所述，$n=1$ 的结果已有基于同余关系的较初等证明。暂且假设 $F$ 为 $\Q$ 的二次扩域，考察 $L(s,\pi,\rho)$ 中可能出现的 $p$ 因子。分两种情形：(1) $p=\mathfrak p_1\mathfrak p_2$ 在 $F$ 中分裂；(2) $p\mathfrak o_F=\mathfrak p$，即 $p$ 在 $F$ 中惰性。

先看 (1)。此时 $G(\Q_p)\simeq\GL_2(F_{\mathfrak p_1})\times\GL_2(F_{\mathfrak p_2})$，其非分歧表示必为 $\pi_1\otimes\pi_2$，其中每个 $\pi_i$ 是 $\GL_2(F_{\mathfrak p_i})$ 的非分歧表示，故对应于 $F_{\mathfrak p_i}^\times$ 的一对非分歧特征 $(\alpha_i,\beta_i)$。更精确地，当由该对特征参数化的整个主系列表示（principal series representation）不可约时，$\pi_i$ 就是该表示；否则它是相应的一维 $\GL_2(F_{\mathfrak p_i})$ 特征。由强逼近，后一情形仅在所考虑的整体表示也是一维时发生。无论哪种情形，由于 $p$ 分裂，局部 $L$ 群 ${}^LG_p$ 只是直积 $\GL_2(\C)\times\GL_2(\C)$，其对应元素为
\[
 g(\pi_p)=\left(
 \begin{pmatrix}\alpha_1(p)&0\\0&\beta_1(p)\end{pmatrix},
 \begin{pmatrix}\alpha_2(p)&0\\0&\beta_2(p)\end{pmatrix}\right).
\]
表示 $\rho$ 就是 $\C^2\otimes\C^2$，因此 $L(s-1,\pi,\rho)$ 的 $p$ 因子为
\begin{equation}
 \begin{aligned}
 &(1-\alpha_1(p)\alpha_2(p)p^{1-s})^{-1}
 (1-\alpha_1(p)\beta_2(p)p^{1-s})^{-1}\\
 &\quad\cdot(1-\beta_1(p)\alpha_2(p)p^{1-s})^{-1}
 (1-\beta_1(p)\beta_2(p)p^{-s})^{-1}.
 \end{aligned}
 \tag{2.2}
\end{equation}
\noindent\emph{译注：所供扫描本中 (2.2) 前三个因子含 $p^{1-s}$，最后一个因子印为 $p^{-s}$；此处按扫描本保留。按前文的统一移位 $L(s-1,\pi,\rho)$，末因子通常也应含 $p^{1-s}$。}

接着看 (2)。此时 $\pi_p$ 是 $\GL_2(F_\mathfrak p)$ 的一个非分歧表示，设对应于 $F_\mathfrak p^\times/\mathfrak o_\mathfrak p^\times$ 的特征 $(\alpha,\beta)$。另一方面，由于 Frobenius $\Phi$ 生成 Galois 群，局部 $L$ 群仍为 $\GL_2(\C)\times\GL_2(\C)$ 与二阶 Galois 群的半直积。直接验证可知，${}^LG_p$ 中的元素 $g(\pi_p)$ 为
\[
 \left(\begin{pmatrix}1&0\\0&1\end{pmatrix},
 \begin{pmatrix}\alpha(p)&0\\0&\beta(p)\end{pmatrix}\right)\Phi.
\]
若 $e_1,e_2$ 为 $\C^2$ 的标准基，则 $g(\pi_p)$ 在 $\rho$ 的空间上的作用如下；回忆 $\Phi$ 交换因子：
\[
 \begin{aligned}
 e_1\otimes e_1&\longmapsto\alpha(p)e_1\otimes e_1,&
 e_1\otimes e_2&\longmapsto\alpha(p)e_2\otimes e_1,\\
 e_2\otimes e_1&\longmapsto\beta(p)e_1\otimes e_2,&
 e_2\otimes e_2&\longmapsto\beta(p)e_2\otimes e_2.
 \end{aligned}
\]
因此，$L(s-1,\pi,\rho)$ 的 $p$ 因子为
\begin{equation}
 (1-\alpha(p)p^{1-s})^{-1}(1-\beta(p)p^{1-s})^{-1}
 (1-\alpha(p)\beta(p)p^{2-2s})^{-1}.
 \tag{2.3}
\end{equation}
当 $n>2$ 时，类似公式会相当复杂。顺便指出，先前提到的 Shimura 未发表结果，恰好是对某些二次域 $F$，其 $p$ 因子具有 (2.2)、(2.3) 的形式。

当 $F=\Q$（$n=1$）时，将 Hasse--Weil $\zeta$ 函数表示为与自守形式相联系的 $L$ 函数的一个推论，是其函数方程和解析延拓。当 $n>1$ 时，这一推论并非自动成立。即使 $n=2$ 已是困难问题，Asai [1] 直到最近才在对 $F$ 施加某种可能并非必要的限制下解决。对 $n>2$ 似乎尚无已知结果。若将 ${}_KS$ 看作 $F$ 上的簇，为其 $\zeta$ 函数建立函数方程相对容易，至少对某些小次数如此；这依赖于 Jacquet 和 Shimura 对 Rankin 思想的推广，例如见 [15] 和 [29]。

\origsection{3}{解析迹公式（The analytical trace formula）}
本节为等式 (2.1) 的右端建立公式，下一节为左端建立公式，第 5 节比较二者。从第 4 节起，我将假设 $p$ 在 $F$ 中分裂，虽然 Langlands 当然没有作这一假设。这个假设大幅简化了论证，同时仍保留若干重要思想。它使我们可以略去整个证明中最复杂、也最引人入胜的一环：使用某个 Bruhat--Tits 建筑（Bruhat--Tits building）中的路径。下面着重介绍其余步骤。Langlands 大概会在自己的报告中谈到这个问题（见 [17]）。

\origsubsection{3.1}{}
等式 (2.1) 的右端
\begin{equation}
 \sum_{\pi\text{ 在 }\mathcal A_0\text{ 中}}
 m(\pi,K)p^{mn/2}\operatorname{trace}\rho(g(\pi_p)^m)
 \tag{3.1}
\end{equation}
可表示为作用于 $\mathcal A_0$ 的某个算子 $R_f$ 的迹，其中 $f\in C_c^\infty(G(\A))$。回忆 $\mathcal A_0$ 为 $G(\A)$ 的不可约、容许、酉表示 $\pi=\pi_\R\otimes\pi^p\otimes\pi_p$ 的离散直和。因此对任意 $f\in C_c^\infty(G(\A)/Z_K)$，有
\[
 \operatorname{trace}R_f=\sum_{\pi\text{ 在 }\mathcal A_0\text{ 中}}
 \operatorname{trace}\pi(f).
\]
当 $f=f_\R f^pf_p$ 时，它又等于
\begin{equation}
 \sum_{\pi\text{ 在 }\mathcal A_0\text{ 中}}
 \operatorname{trace}\pi_\R(f_\R)
 \operatorname{trace}\pi^p(f^p)\operatorname{trace}\pi_p(f_p).
 \tag{3.2}
\end{equation}
回忆 $m(\pi,K)=m(\pi_\R)\dim(\pi^p)^{K^p}$。比较 (3.1)、(3.2)，提示应选择 $f$ 的三个因子，使对 $G(\R)/Z(\R)$ 的任意不可约容许表示 $\pi$，有
\begin{equation}
 \operatorname{trace}\pi(f_\R)=m(\pi);
 \tag{3.3a}
\end{equation}
对 $G(\A_f^p)$ 的任意不可约容许表示 $\pi$，有
\begin{equation}
 \operatorname{trace}\pi(f^p)=\dim\pi^{K^p};
 \tag{3.3b}
\end{equation}
并且
\begin{equation}
 \operatorname{trace}\pi(f_p^{(m)})=
 \begin{cases}
 p^{mn/2}\operatorname{trace}\rho(g(\pi)^m),&\pi\text{ 非分歧},\\
 0,&\text{其他情形}.
 \end{cases}
 \tag{3.3c}
\end{equation}
下面分别构造这三个因子。

$f^p$ 的选择很简单：$f^p=(\operatorname{meas}K^p)^{-1}\operatorname{char}K^p$。$f_p^{(m)}$ 的选择不那么显式，但同样简单：由 Satake 同构（Satake isomorphism）[4]，存在唯一的 $f_p^{(m)}\in\mathcal H(G(\Q_p),G(\Z_p))$ 满足 (3.3c)。$f_\R$ 的构造则更为复杂。

\origsubsection{3.2}{}
若 $G(\R)$ 的表示 $\pi$ 按 $\widetilde G(\R)\simeq\GL_2(\R)^I$\footnote{\textbf{译者注（本次校订）：}按第 1.2 节的定义，$\widetilde G=G/Z$，故此处应为 $\widetilde G(\R)\simeq\operatorname{PGL}_2(\R)^I$；未除去中心的 $G(\R)$ 才同构于 $\GL_2(\R)^I$。} 分解为 $\bigotimes\pi_\iota$，则 $m(\pi)=\prod m(\pi_\iota)$。相应地，若 $f=\prod f_\iota$，则 $\operatorname{trace}\pi(f)=\prod\operatorname{trace}\pi(f_\iota)$。因此，寻找 $f_\R$ 的问题归结为寻找 $f\in C_c^\infty(\operatorname{PGL}_2(\R))$，使对每个不可约容许表示 $\pi$，有
\begin{equation}
 \operatorname{trace}\pi(f)=m(\pi).
 \tag{3.4}
\end{equation}
这可用若干方法做到，例如像 [10] 那样直接由 Paley--Wiener 定理（Paley--Wiener theorem）推出，也可考虑轨道积分（见 [26]）。但 Serge Lang 指出了一个在某种意义下比这两者更初等的论证，下面介绍这一论证。

下文主要依据 Lang 的书 [18, 第 VI 章，第 7 节]。$\operatorname{PGL}_2(\R)$ 的连通分量为 $\operatorname{PSL}_2(\R)$。$\SL_2(\R)$ 中与我们所取 $\GL_2(\R)$ 的极大紧子群对应的极大紧子群为 $\operatorname{SO}(2)$。令 $\varepsilon:\operatorname{SO}(2)\to\C^\times$ 为基本特征（fundamental character）
\[
 \varepsilon\begin{pmatrix}\cos\theta&-\sin\theta\\\sin\theta&\cos\theta\end{pmatrix}
 =\cos\theta+\sqrt{-1}\sin\theta.
\]
$\operatorname{PGL}_2(\R)$ 中唯一一个上同调非平凡的离散系列表示 $\pi_0$，限制到 $\operatorname{SO}(2)$ 后为直和 $\bigoplus\varepsilon^n$（$|n|\ge2$，$n$ 为偶数）。其他离散系列表示的 $\operatorname{SO}(2)$-分解中均不含 $\varepsilon^2$。

回忆所有紧支集、双 $\operatorname{SO}(2)$-有限分布构成的 Hecke 代数，作用于 $\SL_2(\R)$ 的每个容许表示的空间。若 $D$ 满足对所有 $k_1,k_2\in K$，有
\[
 L_{k_1}R_{k_2}D=\varepsilon^m(k_1)\varepsilon^{-n}(k_2)D,
\]
则对这样的每个表示 $(\pi,V)$，算子 $\pi(D)$ 将整个 $V$ 映入 $\varepsilon^m$-特征空间，并消去除 $\varepsilon^n$-特征空间外的一切。在特殊情形 $m=n$ 且 $D$ 仅为在 $K$ 上对 $\varepsilon^{-n}$ 的积分时，$\pi(D)$ 在 $\varepsilon^n$-特征空间上为恒等算子。因此，可选择满足下述条件的 $f\in C_c^\infty(\SL_2(\R))$：
\[
 f(k_1gk_2)=\varepsilon^{-n}(k_1)\varepsilon^{-n}(k_2)f(g),
\]
并使它充分逼近上述 $D$，即选择正的、正规化的、支集充分靠近 $K$ 的 $f$，便可假设 $\pi(f)$ 在此特征空间上也为恒等算子。\footnote{\textbf{译者注（本次校订）：}逼近投影只能先保证作用接近恒等算子。这里随后取 $\pi=\pi_0$、$n=2$，相应特征空间为一维；先取具有所需左右 $K$-型、在该直线上作用非零的 $f$，再除以该作用标量，即可得到恒等作用。此外，非平凡 $K$-型的函数一般取复值，正文的“正”不能理解为该 $f$ 处处非负。}

将此用于 $m=n=2$、$\pi=\pi_0$ 的情形，得到 $f_1\in C_c^\infty(\SL_2(\R))$，使 (a) 对所有 $k_1,k_2\in\operatorname{SO}(2)$、$g\in\SL_2(\R)$，有 $f_1(k_1gk_2)=\varepsilon^{-2}(k_1)\varepsilon^{-2}(k_2)f_1(g)$；(b) $\pi_0(f_1)$ 在 $\varepsilon^2$-特征空间上为恒等算子，在其余 $\operatorname{SO}(2)$-特征空间上为零。由上文关于其他离散系列表示与 $\varepsilon^2$ 的注记，甚至有
\begin{equation}
 \operatorname{trace}\pi(f_1)=\begin{cases}
 1,&\pi=\pi_0,\\
 0,&\pi\text{ 为不同于 }\pi_0\text{ 的任意离散系列表示}.
 \end{cases}
 \tag{3.5}
\end{equation}

回忆 Lang 所称 $f\in C_c^\infty(\SL_2(\R))$ 的 Harish 变换（Harish transform）：
\[
 H_f(a)=|x-x^{-1}|\int_{A\backslash G}f(g^{-1}ag)\,dg,
 \qquad a=\begin{pmatrix}x&0\\0&x^{-1}\end{pmatrix},\quad x\ne\pm1.
\]
这里 $A$ 为所有对角矩阵组成的群。函数 $H_f$ 属于 $C_c^\infty(A)$，甚至属于 $C_c^\infty(A)^W$，其中 $W$ 为 Weyl 群（Weyl group）。实际上，$f\mapsto H_f$ 在双 $\operatorname{SO}(2)$-不变的 $f\in C_c^\infty(\SL_2(\R))$ 的空间与 $C_c^\infty(A)^W$ 之间诱导一个同构 [18, 第 V 章，第 2 节]。因此可选取双 $\operatorname{SO}(2)$-不变的 $f_2$，使 $H_{f_1}=H_{f_2}$。令 $f=f_1-f_2$，则 $H_f=0$。

由于离散系列表示没有被 $\operatorname{SO}(2)$ 固定的非零向量，等式 (3.5) 将 $f_1$ 换为 $f$ 后仍成立。此外，若 $\chi$ 为 $A$ 的任意特征，$\pi_\chi$ 为对应的 $\SL_2(\R)$ 的主系列表示，则
\begin{equation}
 \operatorname{trace}\pi_\chi(f)=\int_AH_f(a)\chi(a)\,da=0
 \tag{3.7}
\end{equation}
[18, 第 VII 章，第 3 节]。由于 $\operatorname{PGL}_2(\R)$ 的任意主系列表示限制到 $\operatorname{PSL}_2(\R)$ 后为某个 $\pi_\chi$，同样结论对 $\operatorname{PGL}_2(\R)$ 的所有主系列表示成立。

最后，若 $\pi$ 为 $\operatorname{PGL}_2(\R)$ 的任意不可约有限维表示，因为它的特征标与某个唯一的离散系列表示的特征标之和为一个主系列表示的特征标，有
\begin{equation}
 \operatorname{trace}\pi(f)=\begin{cases}
 -1,&\pi=\C\text{ 或 }\operatorname{sgn}(\det),\\
 0,&\text{其他有限维 }\pi.
 \end{cases}
 \tag{3.8}
\end{equation}
现在在 $\operatorname{PGL}_2(\R)$ 上定义 $f_\R$：在 $\operatorname{PSL}_2(\R)$ 上取 $-f$，在另一个连通分量上取零。等式 (3.5)、(3.7)、(3.8) 给出 $f=f_\R$ 时的 (3.4)。

\origsubsection{3.3}{}
概括起来，函数 $f=f_\R f^pf_p^{(m)}\in C_c^\infty(G(\A)/Z_K)$ 满足
\[
 \operatorname{trace}(R_f|\mathcal A_0)
 =\sum_{\pi\text{ 在 }\mathcal A_0\text{ 中}}
 m(\pi,K)p^{mn/2}\operatorname{trace}\rho(g(\pi_p)^m).
\]
第 5 节还将用 Selberg 迹公式表示这个迹。为便于稍后比较，这里以相当一般的形式回顾该公式。设 $\mathcal G$ 为任意局部紧群，$\Gamma$ 为离散子群，且 $\Gamma\backslash\mathcal G$ 紧。任意 $\gamma\in\Gamma$ 的共轭类都在 $\mathcal G$ 中闭，这一点可直接验证。因此，对任意 $F\in C_c^\infty(\mathcal G)$，轨道积分
\[
 \int_{\mathcal G_\gamma\backslash\mathcal G}F(g^{-1}\gamma g)\,dg
\]
有定义且有限，其中 $\mathcal G_\gamma$ 为 $\gamma$ 在 $\mathcal G$ 中的中心化子。若 $\Gamma_\gamma=\mathcal G_\gamma\cap\Gamma$，则 $\Gamma_\gamma\backslash\mathcal G_\gamma$ 显然也紧。在某些温和假设下，算子 $R_F$ 在 $L^2(\Gamma\backslash\mathcal G)$ 上为迹类，并且
\[
 \operatorname{trace}R_F=\sum_{\{\gamma\}}
 \operatorname{meas}(\Gamma_\gamma\backslash\mathcal G_\gamma)
 \int_{\mathcal G_\gamma\backslash\mathcal G}F(g^{-1}\gamma g)\,dg,
\]
其中求和遍及 $\Gamma$ 的所有共轭类。第 5 节将把这一公式用于 $\mathcal G=G(\A)/Z_K$、$\Gamma=G(\Q)/(G(\Q)\cap Z_K)$、$F=f$（即刚构造的函数）的情形。

\origsection{4}{代数迹（The algebraic trace）}
如前所述，本节与下一节假定 $p$ 在 $F$ 中完全分裂。

\origsubsection{4.1}{}
$\#{}_KS(\F_{p^m})$ 存在一个与 Selberg 迹公式对等式 (2.1) 右端所给公式极为相似的表达式。

集合 ${}_KS(\overline\F_p)$ 自然地按同源（isogeny）等价关系划分：两个点处于同一类，当且仅当它们参数化的“阿贝尔簇加附加结构”彼此同源。同源类有两种类型：(1) 与系统 $(F',\{\mathfrak q_i\})$ 对应的类，其中 $F'$ 为 $F$ 的全虚二次扩域（totally imaginary quadratic extension），在 $p$ 之上的某个素理想处分裂；$\mathfrak q_i$ 为 $F'$ 的某些理想：若 $\mathfrak p_1,\ldots,\mathfrak p_r$ 为 $F$ 中在 $F'$ 内分裂的素理想，则 $\mathfrak q_i$ 是 $\mathfrak p_i$ 的两个素因子之一。两个系统 $(F'_1,\{\mathfrak q_{1,i}\})$、$(F'_2,\{\mathfrak q_{2,i}\})$ 确定同一同源类，当且仅当 $F'_1=F'_2$ 且这两组理想相等或彼此共轭。(2) 与四元数除代数 $D/F$ 对应的唯一一个同源类，$D$ 恰在以下地方分歧：(a) $B/F$ 分歧的地方；(b) $F$ 中 $p$ 之上的所有素理想；(c) $F$ 的所有实赋值。后一情形可以看作所有不在 $p$ 之上分裂的虚二次扩域 $F'/F$ 的合并。

若 $Y=(F',\{\mathfrak q_i\})$ 或 $D$，将对应于它的 ${}_KS(\overline\F_p)$ 中的点类记为 ${}_KS(Y)$。它在 Frobenius 下稳定，可以相当明确地描述它连同这个作用的结构。关于同源类的参数化及每个类的结构，参见 Milne 的讲座。每个 $Y$ 对应以下数据。

(1) 一个定义在 $\Q$ 上的代数群 $H=H_Y$。若 $Y=(F',\{\mathfrak q_i\})$，则 $H$ 可以说是 $F'$ 的乘法群；若 $Y=D$，则为 $D$ 的乘法群。

(2) 一个定义在 $\Q_p$ 上的代数群 $\overline G_Y=\overline G$。首先，对 $F$ 中每个 $p$ 之上的素理想 $\mathfrak p$，令
\[
 \overline G_\mathfrak p=\begin{cases}
 \text{$F'_\mathfrak p$ 的乘法群}\simeq(F_\mathfrak p^\times)^2,
 &Y=(F',\{\mathfrak q_i\})\text{ 且 }\mathfrak p\text{ 在 }F'\text{ 中分裂},\\
 \text{$D_\mathfrak p$ 的乘法群},&\text{其他情形}.
 \end{cases}
\]
这里“其他情形”指 $Y=(F',\{\mathfrak q_i\})$ 且 $\mathfrak p$ 在 $F'$ 中不分裂，或 $Y=D$。群 $\overline G$ 为 $\prod\overline G_\mathfrak p$。

(3) 陪集空间 $X=\prod X_\mathfrak p$，其中 $X_\mathfrak p=\overline G_\mathfrak p(F_\mathfrak p)/\overline G_\mathfrak p(\mathfrak o_\mathfrak p)$。在每个 $X_\mathfrak p$ 上定义变换 $\Phi_\mathfrak p$：若 $Y=(F',\{\mathfrak q_i\})$、$\mathfrak p$ 在 $F'$ 中分裂，且 $\mathfrak q$ 为位于 $\mathfrak p$ 之上的所选 $\mathfrak q_i$，则 $\Phi_\mathfrak p$ 为乘以 $\mathfrak q$ 的一个一致化元素；否则，为乘以 $D_\mathfrak p$ 的素理想的一个一致化元素。第一种情形若将 $\overline G_\mathfrak p$ 写成 $(F_\mathfrak p^\times)^2$，使第一个因子对应于 $\mathfrak q$，则 $\Phi_\mathfrak p$ 更明确地为乘以 $(p,1)$。令 $X=X_Y=\prod X_\mathfrak p$，$\Phi=\Phi_Y=\prod\Phi_\mathfrak p$。

令 $\overline H(\Q)=H(\Q)/(H(\Q)\cap Z_K)$，并以稍有滥用的记号记
\[
 \overline G(\A_f)=
 \bigl(G(\A_f^p)\times\overline G(\Q_p)\bigr)/(Z_K\cap\A_f^\times).
\]
在每一种情形中，都有：

\begin{originaltheorem}{引理 4.1.1（Lemma 4.1.1）}
(a) 代数群 $Z$ 典范地嵌入 $H$。

(b) 存在 $H(\Q)$ 到 $G(\A_f^p)\times\overline G(\Q_p)$ 的等价嵌入的一个典范类，使 $\overline H(\Q)$ 成为 $\overline G(\A_f)$ 的离散子群。
\end{originaltheorem}

只需解释第二项。显然存在 $H(\Q)$ 到 $G(\A_f^p)$ 及到 $\overline G(\Q_p)$ 的自然嵌入，且它们在相差内自同构的意义下等价。在所有情形中，$H(\R)/Z(\R)$ 都紧，这并非偶然。因此，由 $H(\Q)$ 在 $H(\A)=H(\A_f)\times H(\R)$ 中离散，可知 $\overline H(\Q)$ 在 $H(\A_f)/(Z_K\cap\A_f^\times)$ 中离散，从而在 $\overline G(\A_f)$ 中离散。

Langlands 关于 ${}_KS(Y)$ 结构的主要结果是：其点集合与以下双陪集空间之间存在双射的一个典范类：
\[
 \begin{aligned}
 H(\Q)\backslash\bigl(G(\A_f^p)\times\overline G(\Q_p)\bigr)
 /(K^p\times\overline G(\Z_p))
 &\simeq\overline H(\Q)\backslash\overline G(\A_f)
 /(K^p\times\overline G(\Z_p))\\
 &\simeq H(\Q)\backslash\bigl(G(\A_f^p)\times X\bigr)/K^p.
 \end{aligned}
\]
${}_KS(Y)$ 上的 Frobenius 对应于 $\Phi$ 通过在 $X$ 上的作用而在此陪集空间上诱导的作用。

\origsubsection{4.2}{}
${}_KS(Y)$ 中在 $\F_{p^m}$ 上有理的点的个数，等于 Frobenius 的 $m$ 次幂的不动点数。为表达这一计数，先讨论一个一般性结果。

设 $\mathcal G$ 为任意局部 profinite 群（locally profinite group），$\Gamma$ 为离散子群。则 $\mathcal G$，进而 $C_c^\infty(\mathcal G)$，通过右正则表示作用于 $C^\infty(\Gamma\backslash\mathcal G)$。具体地，对每个 $F\in C_c^\infty(\mathcal G)$、$f\in C^\infty(\Gamma\backslash\mathcal G)$，有
\[
 \begin{aligned}
 R_Ff(x)&=\int_{\mathcal G}F(g)f(xg)\,dg
 =\int_{\mathcal G}F(x^{-1}y)f(y)\,dy\\
 &=\int_{\Gamma\backslash\mathcal G}K_F(x,y)f(y)\,dy,
 \qquad K_F(x,y)=\sum_{\gamma\in\Gamma}F(x^{-1}\gamma y).
 \end{aligned}
\]
若 $K_F(x,x)$ 在 $\Gamma\backslash\mathcal G$ 上紧支，则称算子 $R_F$（有时也称 $F$ 本身）具有形式迹（formal trace），定义为 $\int_{\Gamma\backslash\mathcal G}K_F(x,x)\,dx$。这里使用这一概念，是因为Frobenius 的某个幂的不动点数可以表示为这样的迹，而且有计算公式。

对每个 $\gamma\in\Gamma$，令 $\mathcal G_\gamma$ 为其在 $\mathcal G$ 中的中心化子，$\Gamma_\gamma=\Gamma\cap\mathcal G_\gamma$。若 $F\in C_c^\infty(\mathcal G)$，则在许多情形下轨道积分 $\int_{\mathcal G_\gamma\backslash\mathcal G}F(g^{-1}\gamma g)\,dg$ 有定义；最初等的充分条件是 $F(g^{-1}\gamma g)$ 在 $\mathcal G_\gamma\backslash\mathcal G$ 上紧支，而 $\gamma$ 的共轭类在 $\mathcal G$ 中闭又可保证这一点。

\begin{originaltheorem}{命题 4.2.1（Proposition 4.2.1）}
设 $F\in C_c^\infty(\mathcal G)$。假设满足对某个 $g\in\mathcal G$ 有 $F(g^{-1}\gamma g)\ne0$ 的 $\gamma\in\Gamma$，仅属于 $\Gamma$ 的有限多个共轭类；并假设对每个这样的 $\gamma$，空间 $\Gamma_\gamma\backslash\mathcal G_\gamma$ 紧，且 $\gamma$ 在 $\mathcal G$ 中的共轭类闭。则 $F$ 具有形式迹，其值为
\[
 \sum_{\{\gamma\}}\operatorname{meas}(\Gamma_\gamma\backslash\mathcal G_\gamma)
 \int_{\mathcal G_\gamma\backslash\mathcal G}F(g^{-1}\gamma g)\,dg,
\]
其中求和遍及 $\Gamma$ 的所有共轭类。
\end{originaltheorem}

\begin{proof}
选择足够小的紧开子群 $\mathcal G_0\subset\mathcal G$，使 $\Gamma\cap\mathcal G_0=\{1\}$，且 $F$ 关于 $\mathcal G_0$ 双不变。于是 $\Gamma$ 在 $\mathcal G/\mathcal G_0$ 上自由作用。\footnote{\textbf{译者注（本次校订）：}由 $\Gamma\cap\mathcal G_0=\{1\}$ 只能推出单位陪集的稳定子平凡。自由作用需要对所有 $x$ 都有 $\Gamma\cap x\mathcal G_0x^{-1}=\{1\}$。一般情形可保留这些有限稳定子，并将陪集体积写为 $\operatorname{meas}(\mathcal G_0)/|\Gamma\cap x\mathcal G_0x^{-1}|$；展开积分后仍得到命题中的轨道积分公式。}取代表元集合 $\mathcal X$，使 $\mathcal G$ 为 $\Gamma x\mathcal G_0$（$x\in\mathcal X$）的不交并。由对 $F$ 的一个假设，对每个 $x\in\mathcal X$，和 $\sum_{\gamma\in\Gamma}F(x^{-1}\gamma x)$ 可限于 $\Gamma$ 的有限多个共轭类，故可写成
\[
 \sum_{\{\gamma\}}\sum_{\delta\in\Gamma_\gamma\backslash\Gamma}
 F(x^{-1}\delta^{-1}\gamma\delta x),
\]
其中外层取共轭类代表元，只有有限多个需要计入。由于 $\Gamma_\gamma\backslash\mathcal G_\gamma$ 紧且 $\gamma$ 在 $\mathcal G$ 中的共轭类闭，函数 $F(g^{-1}\gamma g)$ 属于 $C_c^\infty(\Gamma_\gamma\backslash\mathcal G)$；因此上式只在 $\mathcal X$ 的有限子集上非零。故 $F$ 确有形式迹，且等于
\[
 (\operatorname{meas}\mathcal G_0)
 \sum_{x\in\mathcal X}\sum_{\{\gamma\}}
 \sum_{\delta\in\Gamma_\gamma\backslash\Gamma}
 F(x^{-1}\delta^{-1}\gamma\delta x).
\]
实际上，每个和都可限制到有限但任意大的子集。重排后得到
\[
 \begin{aligned}
 &\sum_{\{\gamma\}}\int_{\Gamma_\gamma\backslash\mathcal G}
 F(x^{-1}\gamma x)\,ds\\
 &\quad=\sum_{\{\gamma\}}\operatorname{meas}(\Gamma_\gamma\backslash\mathcal G_\gamma)
 \int_{\mathcal G_\gamma\backslash\mathcal G}F(x^{-1}\gamma x)\,dx.
 \end{aligned}
\]
\end{proof}

\origsubsection{4.3}{}
现在令 $\Gamma=\overline H(\Q)$，$\mathcal G=\overline G(\A_f)$，两者均对应于某个同源类 $Y$。定义
\[
 \begin{aligned}
 F^p&=(\operatorname{meas}K^p)^{-1}\operatorname{char}K^p,\\
 F_\mathfrak p^{(m)}&=(\operatorname{meas}\overline G_\mathfrak p(\mathfrak o_\mathfrak p))^{-1}
 \operatorname{char}\bigl(\Phi_\mathfrak p^m\overline G_\mathfrak p(\mathfrak o_\mathfrak p)\bigr),\\
 F_p^{(m)}&=\prod_{\mathfrak p\mid p}F_\mathfrak p^{(m)},
 \qquad F^{(m)}=F^pF_p^{(m)}.
 \end{aligned}
\]

\begin{originaltheorem}{引理 4.3.1（Lemma 4.3.1）}
对上述 $\Gamma,\mathcal G$ 及 $F=F^{(m)}$（$m\ge1$）的选择，命题 4.2.1 的假设成立。
\end{originaltheorem}

当 $Y=D$ 时显然如此，因为 $\Gamma\backslash\mathcal G$ 实际上紧。因此设 $Y=(F',\{\mathfrak q_i\})$。需要进一步验证的假设，是满足对某个 $g\in\mathcal G$ 有 $F(g^{-1}\gamma g)\ne0$ 的 $\gamma$ 的共轭类数目有限。此时 $H(\Q)=(F')^\times$，$\overline G(\Q_p)=\prod\overline G_\mathfrak p(F_\mathfrak p)$，其中 $\overline G_\mathfrak p$ 为 (a) $(F'_\mathfrak p)^\times$ 或 (b) $D_\mathfrak p^\times$。两种情形中，$F_p^{(m)}(g^{-1}\gamma g)=F_p^{(m)}(\gamma)$，且 $F_p^{(m)}(\gamma)\ne0$ 的充要条件为：(a) 在 $F'$ 分裂的每个 $\mathfrak p_i$ 处，$\gamma$ 在 $\mathfrak q_i$ 处的阶为 $m$，并在其共轭处为单位；(b) 在 $F'$ 不分裂处，$\gamma$ 为 $\mathfrak p_D^m$ 的一个生成元。此外，存在某个 $g\in\overline G(\A)$ 使 $F^p(g^{-1}\gamma g)\ne0$ 的充要条件，是 $\gamma$ 在所有不整除 $p$ 的素理想处为单位。\footnote{\textbf{译者注（本次校订）：}当 $K^p$ 含有额外水平条件时，“为单位”是必要条件，充分性还取决于能否共轭进入 $K^p$。本段证明有限性只使用这一必要条件，因此将“充要条件”改为“必要条件”即可。}因此若 $F^{(m)}(g^{-1}\gamma g)\ne0$，$\gamma$ 在 $F'$ 所有素理想处的阶均已确定。任意两个这样的 $\gamma$ 相差一个单位，而 $F'$ 的单位模去 $Z_K$ 后为有限集合。

函数 $F^{(m)}$ 的重要性在于：

\begin{originaltheorem}{命题 4.3.2（Proposition 4.3.2）}
$F^{(m)}$ 的形式迹等于 Frobenius 的 $m$ 次幂在 ${}_KS(Y)$ 上的不动点个数。
\end{originaltheorem}

这个简单练习留给读者。

\origsection{5}{比较（The comparison）}
\origsubsection{5.1}{}
由第 3 节可知，等式 (2.1) 的右端为
\begin{equation}
 \sum_{\{\gamma\}}\operatorname{meas}(\Gamma_\gamma\backslash\mathcal G_\gamma)
 \int_{\mathcal G_\gamma\backslash\mathcal G}F(g^{-1}\gamma g)\,dg,
 \tag{5.1}
\end{equation}
其中 $\Gamma=G(\Q)/(G(\Q)\cap Z_K)$，$\mathcal G=G(\A)/Z_K$，$F=f_\R f^pf_p^{(m)}$。

令 $\varphi:G(\A)\to\mathcal G$ 为典范投影。容易看出 $\varphi(G_\gamma(\A))\backslash\mathcal G_\gamma$ 总是紧，因而上述公式中可用 $\varphi(G_\gamma(\A))$ 替换 $\mathcal G_\gamma$。再注意
\[
 \varphi(G_\gamma(\A))\backslash\mathcal G\simeq G_\gamma(\A)\backslash G(\A).
\]
利用 $G(\A)$ 的分解，每一项可写成以下四个因子的乘积：
\begin{equation}
 \operatorname{meas}(G_\gamma(\Q)Z_K\backslash G_\gamma(\A)),
 \tag{5.2a}
\end{equation}
\begin{equation}
 \int_{G_\gamma(\A_f^p)\backslash G(\A_f^p)}f^p(g^{-1}\gamma^pg)\,dg,
 \tag{5.2b}
\end{equation}
\begin{equation}
 \int_{G_\gamma(\Q_p)\backslash G(\Q_p)}f_p^{(m)}(g^{-1}\gamma_pg)\,dg,
 \tag{5.2c}
\end{equation}
\begin{equation}
 \int_{G_\gamma(\R)\backslash G(\R)}f_\R(g^{-1}\gamma_\R g)\,dg.
 \tag{5.2d}
\end{equation}
由第 4 节可知，$\#{}_KS(\F_{p^m})$ 可表示为
\begin{equation}
 \sum_Y\sum_{\{\gamma\}}\operatorname{meas}(\Gamma_\gamma\backslash\mathcal G_\gamma)
 \int_{\mathcal G_\gamma\backslash\mathcal G}F^{(m)}(g^{-1}\gamma g)\,dg,
 \tag{5.3}
\end{equation}
其中 $Y$ 遍及所有系统 $(F',\{\mathfrak q_i\})$（将一个系统与其共轭等同），以及唯一的 $D$；$\gamma$ 遍及 $H_Y(\Q)$ 的所有共轭类；$\mathcal G=\overline G_Y(\A_f)=G(\A_f^p)\times\overline G_Y(\Q_p)$；$F^{(m)}=F^pF_p^{(m)}$，同第 4 节末。与前面一样，每一项都分解为以下三个因子的乘积：
\begin{equation}
 \operatorname{meas}\bigl(H_\gamma(\Q)(Z_K\cap\A_f^\times)
 \backslash G_\gamma(\A_f^p)\overline G_\gamma(\Q_p)\bigr),
 \tag{5.4a}
\end{equation}
\begin{equation}
 \int_{G_\gamma(\A_f^p)\backslash G(\A_f^p)}F^p(g^{-1}\gamma^pg)\,dg,
 \tag{5.4b}
\end{equation}
\begin{equation}
 \int_{\overline G_\gamma(\Q_p)\backslash\overline G(\Q_p)}
 F_p^{(m)}(g^{-1}\gamma_pg)\,dg.
 \tag{5.4c}
\end{equation}
主要定理的证明归结为测度与轨道积分的比较。

\origsubsection{5.2}{}
先比较两边使用的测度。暂令 $k$ 为任意局部域，$\psi$ 为其加法特征。若 $H$ 为 $k$ 上的有限维半单代数，$\Tr$ 为其约化迹（reduced trace），则 $\psi\circ\Tr$ 为 $H$ 的一个加法特征。$H$ 上存在唯一测度 $dx$，使其在由此特征确定的 Fourier 变换下自对偶；继而在 $H^\times$ 上得到不变测度 $d^\times x=dx/|x|_H$。

若 $k$ 为整体域，$\psi=\prod\psi_v$ 为 $\A_k/k$ 的特征，$H$ 为 $k$ 上的半单代数，则由此在所有局部群 $H_v^\times$ 及 $H^\times(\A_k)$ 上得到测度，称为由 $\psi$ 确定的 Tamagawa 测度（见 [16] 第 15 节）。依次将此构造应用于域 $F$ 本身、$F$ 的二次扩域和 $F$ 上的四元数代数（包括 $M_2(F)$）。以下总是假定使用这样选择的测度。一个经典结果是，整体测度
\[
 \operatorname{meas}\bigl(H^\times(\A)/(\A_F^\times\cdot H^\times(F))\bigr)
\]
与四元数代数 $H$ 无关（[16] 第 16 节，或 Gelbart 与 Jacquet 的讲座 [12]）。还需要局部测度的比较。

\begin{originaltheorem}{引理 5.2.1（Lemma 5.2.1）}
设 $k$ 为任意非 Archimedes 局部域，$H$ 为 $k$ 上唯一的四元数除代数。则
\[
 \operatorname{meas}\mathfrak o_H^\times
 =(q-1)^{-1}\operatorname{meas}\GL_2(\mathfrak o_k).
\]
\end{originaltheorem}

\begin{proof}
改变 $\psi$ 只会将两个测度乘以相同的常数；选择 $\psi$ 的导子（conductor）为 $\mathfrak o_k$。令 $f$ 为 $M_2(\mathfrak o_k)$ 的示性函数，则
\[
 \begin{aligned}
 \widehat f(x)&=\int_{M_2(k)}f(y)\psi(\Tr(xy))\,dy\\
 &=\operatorname{meas}M_2(\mathfrak o_k)
 \begin{cases}1,&x\in M_2(\mathfrak o_k),\\0,&x\notin M_2(\mathfrak o_k),\end{cases}\\
 &=(\operatorname{meas}M_2(\mathfrak o_k))f(x).
 \end{aligned}
\]
因此 $\operatorname{meas}M_2(\mathfrak o_k)=1$。群 $\GL_2(\mathfrak o_k)$ 在 $M_2(\mathfrak o_k)$ 中开，而且 $x\in M_2(\mathfrak o_k)$ 属于 $\GL_2(\mathfrak o_k)$ 当且仅当 $x\bmod\mathfrak p_k$ 非奇异，故
\[
 \operatorname{meas}\GL_2(\mathfrak o_k)
 =\frac{\#\GL_2(\F_q)}{\#M_2(\F_q)}
 =(1-1/q)^2(1+1/q).
\]
现在令 $f$ 为 $\mathfrak o_H$ 的示性函数，则
\[
 \widehat f(x)=\int_{\mathfrak o_H}\psi(\Tr(xy))\,dy
 =\operatorname{meas}\mathfrak o_H
 \begin{cases}1,&x\in\mathfrak p_H^{-1},\\0,&x\notin\mathfrak p_H^{-1};\end{cases}
\]
即 $\widehat f=(\operatorname{meas}\mathfrak o_H)\operatorname{char}(\mathfrak p_H^{-1})$。其 Fourier 变换进一步为
\[
 \widehat{\widehat f}
 =(\operatorname{meas}\mathfrak o_H)(\operatorname{meas}\mathfrak p_H^{-1})f.
\]
由于 $\operatorname{meas}\mathfrak p_H^{-1}=q^2\operatorname{meas}\mathfrak o_H$，有 $\operatorname{meas}\mathfrak o_H=q^{-1}$。与上面同理，
\[
 \operatorname{meas}\mathfrak o_H^\times=(1-1/q)(1+1/q)(1/q).
\]
引理由此得证。
\end{proof}

也可类似地比较 $\GL_2(\R)$ 与 $\mathbb H^\times$ 上的测度，不过此处并不需要。

\origsubsection{5.2（续）}{}
再暂令 $k$ 为特征为 $0$ 的任意局部域，$\mathcal G$ 为 $\GL_2(k)$、$k$ 上四元数除代数的乘法群 $H^\times(k)$、$\operatorname{PGL}_2(k)$ 或 $H^\times/k^\times$。对任意 $x\in\mathcal G$，若其在 $k$ 的代数闭包上的特征值为 $\alpha,\beta$，定义 $D(x)=|\alpha-\beta|^2/|\alpha\beta|$。若 $T$ 为 $\mathcal G$ 的任意环面，则对正则 $t\in T$ 和 $f\in C_c^\infty(\mathcal G)$，定义
\[
 O^T(f,t)=\int_{T\backslash\mathcal G}f(g^{-1}tg)\,dg.
\]
函数 $F(t)=O^T(f,t)$ 满足
\begin{equation}
 F\text{ 在 }T\text{ 的正则元素上光滑};\tag{5.5a}
\end{equation}
\begin{equation}
 D^{1/2}F\text{ 局部有界};\tag{5.5b}
\end{equation}
\begin{equation}
 F\text{ 在 }T\text{ 上紧支}.\tag{5.5c}
\end{equation}
若 $\pi$ 为 $\mathcal G$ 的不可约容许表示，则其特征标 $\operatorname{ch}_\pi$ 也满足 (5.5a)--(5.5b)。对任意 $f\in C_c^\infty(\mathcal G)$，有
\[
 \int f(g)\,dg=\frac12\sum_{\{T\}}\int_TD(t)O^T(f,t)\,dt,
\]
其中求和遍及 $\mathcal G$ 中所有环面的共轭类。若 $\pi$ 不可约容许，则
\begin{equation}
 \operatorname{trace}\pi(f)=\frac12\sum_{\{T\}}
 \int_TD(t)\operatorname{ch}_\pi(t)O^T(f,t)\,dt.
 \tag{5.6}
\end{equation}
关于这些内容见 [16] 第 7 节；轨道积分的更多内容见 [12] 和 [26]。

\begin{originaltheorem}{引理 5.2.1（Lemma 5.2.1）}
设 $\mathcal G=H^\times$，$F$ 为 $\mathcal G^{\mathrm{reg}}$ 上共轭不变的函数，其在每个环面 $T$ 上的限制均满足 (5.5a)--(5.5c)。若对 $\mathcal G$ 的所有不可约容许表示 $\sigma$，都有
\[
 \frac12\sum_{\{T\}}\int_TD(t)\operatorname{ch}_\sigma(t)F(t)\,dt=0,
\]
则 $F=0$。
\end{originaltheorem}

固定环面 $T$。$\mathcal G$ 中与 $T$ 的正则元素共轭的元素构成开集。若 $f\in C_c^\infty(T^{\mathrm{reg}})$，则存在支集包含于此开集的 $h\in C_c^\infty(\mathcal G^{\mathrm{reg}})$，使 $O^T(h)=f$。于是
\[
 \int_{\mathcal G}h(g)F(g)\,dg
 =\int_TD(t)O^T(h,t)F(t)\,dt
 =\int_TD(t)f(t)F(t)\,dt.
\]
但不可约表示的特征标在 $\mathcal G$ 的中心函数空间中完备，故 $f=\sum f_\sigma\operatorname{ch}_\sigma$；当 $k$ 非 Archimedes 时，可将此问题归结为有限群的问题。因此根据假设，上述积分为零。换言之，$DF$ 对任意 $f\in C_c^\infty(T^{\mathrm{reg}})$ 的积分为零，故 $F$ 本身为零。

可以且将要把 $H^\times$ 的共轭类与 $\GL_2(k)$ 的椭圆及中心共轭类等同。由 [16] 第 15 节（亦见 [12]），以下两类表示之间存在双射 $\sigma\leftrightarrow\pi$：(a) $H^\times$ 的不可约光滑表示 $\sigma$，因此为有限维；(b) $\GL_2(k)$ 的模中心 $Z$ 平方可积的不可约容许表示 $\pi$。此对应满足，对所有正则椭圆元素 $t$，
\begin{equation}
 \operatorname{ch}_\sigma(t)=-\operatorname{ch}_\pi(t).
 \tag{5.7}
\end{equation}
$H^\times/k^\times$ 的表示对应于 $\operatorname{PGL}_2(k)$ 的表示。当 $k=\R$ 时，$H^\times$ 的平凡表示对应于 $\operatorname{PGL}_2$ 的唯一离散系列表示 $\pi_0$。当 $k$ 非 Archimedes 时，$H^\times$ 的一维表示对应于 Steinberg 表示（Steinberg representation），亦称特殊表示；其余表示对应于 $\GL_2$ 的绝对尖点表示（absolutely cuspidal representation）。

\begin{originaltheorem}{命题 5.2.2（Proposition 5.2.2）}
令 $f_\R\in C_c^\infty(\operatorname{PGL}_2(\R))$ 如第 3 节所构造。对半单元素 $x\in\operatorname{PGL}_2(\R)=\mathcal G$，有
\[
 \int_{\mathcal G_x\backslash\mathcal G}f_\R(g^{-1}xg)\,dg
 =\begin{cases}
 0,&x\text{ 双曲},\\
 (\operatorname{meas}\C^\times/\R^\times)^{-1},&x\text{ 椭圆},\\
 -(\operatorname{meas}\mathbb H^\times/\R^\times)^{-1},&x=1.
 \end{cases}
\]
\end{originaltheorem}
\noindent\emph{译注：所供扫描本此命题的积分中，括号 $(g^{-1}xg)$ 前未见函数名；这里依据命题的上下文补出 $f_\R$。本页引理号 5.2.1 与上一页重复，亦保留扫描本编号。}

\begin{proof}
双曲 $x$ 的情形已包含于 $f_\R$ 的构造。回忆
\[
 \operatorname{trace}\pi(f_\R)=\begin{cases}
 1,&\pi=1\text{ 或 }\operatorname{sgn}(\det),\\
 -1,&\pi=\pi_0,\\0,&\text{其他情形}.
 \end{cases}
\]
现在考虑 $\mathbb H^\times/\R^\times$ 上值为 $(\operatorname{meas}\mathbb H^\times/\R^\times)^{-1}$ 的常值函数 $h_\R$。它满足 $\operatorname{trace}\sigma(h_\R)=1$（$\sigma=1$），其余为零。因此，由 (5.7) 及随后注记，$\pi,\sigma$ 对应时，$\operatorname{trace}\pi(f_\R)=-\operatorname{trace}\sigma(h_\R)$。与 (5.6) 联合得到
\[
 \begin{aligned}
 &\int_{\C^\times/\R^\times}D(t)\operatorname{ch}_\pi(t)\,dt
 \int_{\C^\times\backslash\operatorname{PGL}_2(\R)}f_\R(g^{-1}tg)\,dg\\
 &\qquad=-\int_{\C^\times/\R^\times}D(t)\operatorname{ch}_\sigma(t)\,dt
 \int_{\C^\times\backslash\mathbb H^\times}h_\R(g^{-1}tg)\,dg.
 \end{aligned}
\]
由 (5.7) 和引理 5.2.1（此处其实只需 $\C^\times/\R^\times$ 上的 Fourier 反演（Fourier inversion）），对所有 $t\ne\pm1\in\C^\times/\R^\times$，有
\[
 \begin{aligned}
 \int_{\C^\times\backslash\operatorname{PGL}_2(\R)}f_\R(g^{-1}tg)\,dg
 &=\int_{\C^\times\backslash\mathbb H^\times}h_\R(g^{-1}tg)\,dg\\
 &=(\operatorname{meas}\C^\times/\R^\times)^{-1}.
 \end{aligned}
\]
由于 $f_\R$ 的双曲轨道积分为零，$f_\R$ 与 $-h_\R$ 在 [12] 中“匹配轨道积分”（matching orbital integrals）一节的意义下相对应。故
\[
 f_\R(1)=-h_\R(1)=-(\operatorname{meas}\mathbb H^\times/\R^\times)^{-1}.
\]
\end{proof}

\origsubsection{5.3}{}
现在设 $k$ 为非 Archimedes 局部域，暂令 $G=\GL_2(k)$，$K=\GL_2(\mathfrak o_k)$。对每个 $m\ge0$，令 $f^{(m)}$ 为 Hecke 代数 $\mathcal H(G,K)$ 中由 Satake 同构确定的唯一函数，使当 $\pi=\pi(\alpha,\beta)$ 是由 $k^\times$ 的特征 $x\mapsto(\alpha^{\operatorname{ord}(x)},\beta^{\operatorname{ord}(x)})$ 对应的非分歧主系列表示时，有
\[
 \operatorname{trace}\pi(f^{(m)})=q^{m/2}(\alpha^m+\beta^m).
\]
$f^{(m)}$ 可相当显式地给出。对每个 $i,j\in\Z$，回忆 Hecke 算子
\[
 T(\mathfrak p^i,\mathfrak p^j)
 =(\operatorname{meas}K)^{-1}\operatorname{char}
 \left(K\begin{pmatrix}\eta^i&0\\0&\eta^j\end{pmatrix}K\right),
\]
其中 $\eta$ 为 $\mathfrak p$ 的一个生成元。对 $m\ge0$，还令\footnote{\textbf{译者注（本次校订）：}求和应对无序对 $\{i,j\}$ 各取一次，例如加上 $i\ge j$。由于 $T(\mathfrak p^i,\mathfrak p^j)=T(\mathfrak p^j,\mathfrak p^i)$，遍历全部有序对会重复计数；在 $m=1$ 时便会将 $T(\mathfrak p)$ 错计为两倍，无法满足下述递推关系。} $T(\mathfrak p^m)=\sum T(\mathfrak p^i,\mathfrak p^j)$，其中 $i,j\ge0$，$i+j=m$。这些就是通常的 Hecke 算子，它们满足
\[
 T(\mathfrak p)T(\mathfrak p^m)
 =T(\mathfrak p^{m+1})+qT(\mathfrak p,\mathfrak p)T(\mathfrak p^{m-1}),
\]
形式地写成
\begin{equation}
 (I-T(\mathfrak p)X+qT(\mathfrak p,\mathfrak p)X^2)^{-1}
 =\sum_mT(\mathfrak p^m)X^m.
 \tag{5.7}
\end{equation}

\begin{originaltheorem}{引理 5.3.1（Lemma 5.3.1，Ihara）}
有
\[
 f^{(0)}=2I,\qquad f^{(1)}=T(\mathfrak p),\qquad
 f^{(m)}=T(\mathfrak p^m)-qT(\mathfrak p,\mathfrak p)T(\mathfrak p^{m-2}),\quad m\ge2.
\]
\end{originaltheorem}

\begin{proof}
$m=0$ 显然。众所周知，当 $\pi=\pi(\alpha,\beta)$ 时，$\operatorname{trace}\pi(T(\mathfrak p))=q^{1/2}(\alpha+\beta)$，得到 $m=1$。形式地，$T(\mathfrak p)\leftrightarrow q^{1/2}(\alpha+\beta)$，且容易证明 $T(\mathfrak p,\mathfrak p)\leftrightarrow\alpha\beta$。于是
\[
 \begin{aligned}
 &(1-q^{1/2}\alpha X)^{-1}+(1-q^{1/2}\beta X)^{-1}\\
 &\quad=\sum_{m=0}^\infty q^{m/2}(\alpha^m+\beta^m)X^m
 \longleftrightarrow\sum_mf^{(m)}X^m.
 \end{aligned}
\]
而此表达式又等于
\[
 \frac{(1-q^{1/2}\alpha X)+(1-q^{1/2}\beta X)}
 {(1-q^{1/2}\alpha X)(1-q^{1/2}\beta X)}
 \longleftrightarrow
 \frac{2-T(\mathfrak p)X}{1-T(\mathfrak p)X+qT(\mathfrak p,\mathfrak p)X^2}.
\]
结合生成函数等式 (5.7)，比较系数即得结论。
\end{proof}

\begin{originaltheorem}{推论 5.3.2（Corollary 5.3.2）}
若 $x=\begin{pmatrix}\eta^l&0\\0&\eta^l\end{pmatrix}$，则\footnote{\textbf{译者注（本次校订）：}推论 5.3.2 的公式适用于 $m\ge1$；$m=0$ 时应使用 $f^{(0)}=2I$，故 $f^{(0)}(1)=2/\operatorname{meas}K$。}
\[
 f^{(m)}(x)=\begin{cases}
 -(\operatorname{meas}K)^{-1}(q-1),&m=2l,\\0,&\text{其他情形}.
 \end{cases}
\]
\end{originaltheorem}



仍令 $H$ 为 $k$ 上的四元数除代数。对每个 $m\ge0$，令\footnote{\textbf{译者注（本次校订）：}归一化分母应为乘法群的单位子群体积 $\operatorname{meas}(\mathfrak o_H^\times)$。这使每个赋值壳层上的积分等于 $1$，也与引理 5.2.1 及下文 (c) 的等式相符。若使用此前的加法体积 $\operatorname{meas}(\mathfrak o_H)=q^{-1}$，这些等式便不成立。}
\[
 h^{(m)}=(\operatorname{meas}\mathfrak o_H)^{-1}
 \operatorname{char}(\mathfrak p_H^m\setminus\mathfrak p_H^{m+1}).
\]

\begin{originaltheorem}{命题 5.3.3（Proposition 5.3.3）}
设 $x$ 为 $G$ 的半单元素，$m\ge1$。

(a) 若 $x$ 双曲，则
\[
 \int_{G_x\backslash G}f^{(m)}(g^{-1}xg)\,dg
 =\begin{cases}
 1,&x=\diag(\eta^m,1)\text{ 或 }\diag(1,\eta^m)\pmod{K\cap A},\\
 0,&\text{其他情形}.
 \end{cases}
\]

(b) 若 $x$ 椭圆，则
\[
 \int_{G_x\backslash G}f^{(m)}(g^{-1}xg)\,dg
 =\int_{H_x^\times\backslash H^\times}h^{(m)}(g^{-1}xg)\,dg.
\]

(c) 若 $x$ 为中心元素，则 $f^{(m)}(x)=-h^{(m)}(x)$。
\end{originaltheorem}

\begin{proof}
回忆 Satake 同构与轨道积分之间的关系，例如见 [4]：对 $f\in\mathcal H(G,K)$，有
\[
 \widetilde f(a)=D(a)^{1/2}\int_{A\backslash G}f(g^{-1}ag)\,dg,
 \qquad\operatorname{trace}\pi_\chi(f)=\int_A\widetilde f(a)\chi(a)\,da.
\]
由 $f^{(m)}$ 的定义，
\[
 \widetilde f^{(m)}(a)=\begin{cases}
 q^{m/2},&a=\diag(\eta^m,1)\text{ 或 }\diag(1,\eta^m),\\
 0,&\text{其他情形}.
 \end{cases}
\]
由此得到 (a)。(c) 直接来自 5.3.2 和 5.2.1。只有 (b) 较困难。首先，令 $\pi$ 遍及 $\GL_2$ 的本质平方可积（essentially square-integrable）不可约表示，考虑 (5.6)：
\[
 \begin{aligned}
 \operatorname{trace}\pi(f^{(m)})={}&\frac12\int_A D(a)\operatorname{ch}_\pi(a)\,da
 \int_{A\backslash G}f^{(m)}(g^{-1}ag)\,dg\\
 &+\frac12\sum_{\{T\}}\int_TD(t)\operatorname{ch}_\pi(t)\,dt
 \int_{T\backslash G}f^{(m)}(g^{-1}tg)\,dg,
 \end{aligned}
\]
其中 $T$ 遍及所有各向异性环面（anisotropic torus）。类似地，令 $\sigma$ 遍及 $H^\times$ 的所有不可约容许表示，考虑 (5.6)：
\[
 \operatorname{trace}\sigma(h^{(m)})
 =\frac12\sum_{\{T\}}\int_TD(t)\operatorname{ch}_\sigma(t)\,dt
 \int_{T\backslash H^\times}h^{(m)}(g^{-1}tg)\,dg.
\]
应用 5.2.1 与 (5.7)，只需证明当 $\sigma,\pi$ 对应时，
\[
 -\operatorname{trace}\pi(f^{(m)})
 +\frac12\int_AD(a)\operatorname{ch}_\pi(a)\,da
 \int_{A\backslash G}f^{(m)}(g^{-1}ag)\,dg
 =\operatorname{trace}\sigma(h^{(m)}).
\]
为此：(1) 这里出现的 $\pi$ 都没有非零的 $K$-不动向量，因此始终有 $\operatorname{trace}\pi(f^{(m)})=0$。(2) 若 $\pi$ 为绝对尖点的，则除非 $a\in A(\mathfrak o_k)$，否则 $\operatorname{ch}_\pi(a)=0$；此事首先见于 [19] 引理 6.4，亦见 [8] 和 [5]。根据已证明的 (a)，$O^A(f)$ 在 $A(\mathfrak o)$ 上为零，因此整个左端消失。此时 $\sigma$ 不是一维的，故在 $\mathfrak o_H^\times$ 上不平凡，整个右端也消失。(3) 若 $\pi$ 为特殊表示，对应于 $H^\times$ 的特征 $\sigma=\chi\circ N_H$，则一方面 $\sigma(h^{(m)})=\chi(\eta^m)$\footnote{\textbf{译者注（本次校订）：}该值适用于 $\chi$ 非分歧的情形。若 $\chi$ 在 $\mathfrak o_k^\times$ 上非平凡，单位群上的积分为零，等式两端都应为零。}；另一方面，再次由 (a)，
\[
 \frac12\int_AD(a)\operatorname{ch}_\pi(a)\,da
 \int_{A\backslash G}f^{(m)}(g^{-1}ag)\,dg=\chi(\eta^m).
\]
\end{proof}

\origsubsection{5.4}{回到先前的记号}
最后建立 (5.1) 与 (5.3) 非零项之间的双射，并证明对应项相等。

由 5.2.2，除非 $\gamma_\R$ 椭圆或为中心元素，否则 (5.2) 中 $f_\R$ 的轨道积分为零。因此今后假定属于这两种情形。

首先设 $\gamma\in G(\Q)$ 为中心元素。它在 (5.1) 中对应的项为以下四者的乘积：
\begin{equation}
 \operatorname{meas}(G(\Q)Z_K\backslash G(\A)),\tag{5.8a}
\end{equation}
\begin{equation}
 f_p^{(m)}(\gamma_p)=(-1)^n(\operatorname{meas}D^\times(\Z_p))^{-1},\tag{5.8b}
\end{equation}
\begin{equation}
 f^p(\gamma^p),\tag{5.8c}
\end{equation}
\begin{equation}
 f_\R(\gamma_\R)=(-1)^n(\operatorname{meas}D^\times(\R)/\R^\times)^{-1}.\tag{5.8d}
\end{equation}
在 (5.3) 中，由 $Y=D$ 编号的项中恰有一项与 $\gamma$ 对应。它是以下三者的乘积：
\begin{equation}
 \operatorname{meas}\bigl(D^\times(\Q)(Z_K\cap\A_f^\times)\backslash D^\times(\A_f)\bigr),\tag{5.9a}
\end{equation}
\begin{equation}
 F_p^{(m)}(\gamma_p)=(\operatorname{meas}D^\times(\Z_p))^{-1},\tag{5.9b}
\end{equation}
\begin{equation}
 F^p(\gamma^p).\tag{5.9c}
\end{equation}
这两个乘积相等，因为 (5.8c)、(5.9c) 显然相等，而
\[
 \begin{aligned}
 \operatorname{meas}(G(\Q)Z_K\backslash G(\A))
 &=\operatorname{meas}(D^\times(\Q)Z_K\backslash D^\times(\A))\\
 &=\operatorname{meas}\bigl(D^\times(\Q)(Z_K\cap\A_f^\times)
 \backslash D^\times(\A_f)\bigr)\operatorname{meas}(Z(\R)\backslash D^\times(\R)).
 \end{aligned}
\]

现在设 $\gamma\in G(\Q)$ 的实分量 $\gamma_\R$ 椭圆。则 $F'=F(\gamma)$ 为虚二次扩域。若该扩域在 $F$ 中 $p$ 之上的任何素理想处都不分裂，令 $Y=Y_\gamma=D$。否则，$Y_\gamma$ 为某个 $(F',\{\mathfrak q_i\})$，需要指定这些 $\mathfrak q_i$。由 $f_p^{(m)}$ 是各 $f_\mathfrak p^{(m)}$ 的乘积的定义，以及 5.3.3，可知如果 $\gamma$ 对应的 (5.2c) 非零，则当 $\mathfrak p$ 在 $F'$ 中分裂时，$\gamma_\mathfrak p$ 在 $\GL_2(F_\mathfrak p)$ 中共轭于 $\diag(\eta^m,1)$。换言之，$\gamma_\mathfrak p$ 必须在 $F'$ 中位于 $\mathfrak p$ 之上的某个素理想 $\mathfrak q_\gamma$ 处的阶为 $m$，在另一个处为单位。令 $Y_\gamma=(F',\{\mathfrak q_\gamma\})$。

无论哪种情形，元素 $\gamma$ 显然给出 $H_{Y_\gamma}(\Q)$ 的元素。由 5.2.2 与 5.3.3 容易推出，(5.2) 与 (5.4) 的对应项相等。

\origsection{6}{补充注记（Supplementary remarks）}
\origsubsection{6.1}{}
当 $F=\Q$ 时，主要定理来自 Hecke 算子与 Frobenius 之间的一个同余关系。虽然这样的关系很可能在很一般的情形下成立，但它不能给出 $\zeta$ 函数的公式。下面稍详细地说明。

首先，整系数 Hecke 算子总是在概形 ${}_KS$ 上定义代数对应（algebraic correspondence）。当 $F=\Q$ 时，同余关系说，模 $p$ 有
\[
 T(p)=\Phi+\Phi^*T(p,p),
\]
其中 $\Phi$ 为 Frobenius，$\Phi^*$ 为其转置（参见 [28] 第 7.4 节）。正如 Langlands 向我指出，等价的表述是：在代数对应环中，$\Phi$ 是多项式 $X^2-T(p)X+pT(p,p)$ 的一个根。一般地，可把多项式 $\det(X-\rho(g))$ 视为 ${}^LG$ 中半单元素 $g$ 的函数。通过 Satake 同构，它对应于一个系数为 Hecke 算子的多项式。这些系数实际上是整的，看起来证明 $\Phi$ 是其根并不十分困难。具体而言，这立即意味着作用于 $l$ 进上同调的 Frobenius 的根包含于这个多项式的根中，或者说 $Z_p(X,{}_KS)$ 的因子包含于
\[
 \prod_{\pi\text{ 在 }\mathcal A_0\text{ 中}}
 \det(I-\rho(g(\pi_p)))^{m(\pi,K^p)}
\]
的因子之中。\footnote{\textbf{译者注（本次校订）：}这里应保留多项式变量，并与第 2.2 节的移位一致：相应局部行列式为 $\det(I-p^{n/2}\rho(g(\pi_p))X)$，取 $X=p^{-s}$ 后才是 $\zeta$ 函数的局部因子。正文所写 $\det(I-\rho(g(\pi_p)))$ 既缺少变量，也缺少这一归一化因子。}当 $F=\Q$ 时，利用一个我一直觉得有些神秘的进一步关系，可得到这些根实际上完全相同；但这第二步似乎不能推广。我指的是 [28] 的 7.10(2)，亦参见 [24] 第 333--336 页 Piatetskii-Shapiro 给出的表示论类似论证。Langlands 指出，在某种意义下这种困难有充分理由：当出现 $L$ 不可区分性问题时，Frobenius 的根与 Hecke 多项式的根可能在某种意义下只能部分重合。

顺便说，Shimura 在前述情形证明主要定理时，也使用了 Selberg 迹公式。似乎 Ihara 是第一个将迹公式应用于这类问题的人：在 [14] 中，除去某些技术问题，他证明与 $\SL_2(\Z)\backslash\mathcal H$ 上某些层相联系的 Hasse--Weil $\zeta$ 函数是 Hecke $L$ 函数的乘积。他与 Shimura 都使用可称为整体的表述。第一次涉及表示论与局部轨道积分的应用，是 Langlands 的 Antwerp 报告 [19]；这项早期工作似乎既在 Langlands 后来的发展中发挥作用，也影响了 Drinfeld 的工作，注意 [9] 对此的注记。

\origsubsection{6.2}{}
主要定理的证明只需少量额外工作，就可推广到非平凡层，并作一个精细化，按照 $\mathcal A_0$ 的直和分解分解 $\zeta$ 函数。

先概述如何处理层。设 $\xi:G\to\GL_n$ 为 $G$ 的有理表示（即在 $\Q$ 上有理），并在 $N^1_{F/\Q}$ 上平凡。这里 $N^1_{F/\Q}$ 是从 $F^\times$ 到 $\Q^\times$ 的范数映射的核，作为代数群典范地以标量嵌入 $G$ 的中心。可为 $\xi$ 配上 ${}_KS(\C)$ 上的局部常值 $\Q$-层 $E_\xi$，以及 ${}_KS$ 上的 $l$ 进层 $E_\xi(\Q_l)$，如 [19] 的一个特殊情形，并对每个“好”素数 $p$，考虑一个 $\zeta$ 函数的 $p$ 因子
\[
 \zeta_p(s,{}_KS,\xi)=\prod_{i=0}^{2n}
 \det\bigl(I-(\Phi|H^i({}_KS,E_\xi(\Q_l)))p^{-s}\bigr)^{(-1)^i}.
\]
\noindent\emph{译注：所供扫描本此处指数为 $(-1)^i$；按下文 $\log\zeta_p$ 的正号及通常的 $\zeta$ 函数约定，常见写法为 $(-1)^{i+1}$。这里保留扫描本所示指数。}

由 Lefschetz 迹公式（Lefschetz trace formula），它满足
\[
 \log\zeta_p=\sum_{m=1}^\infty\frac1{mp^{ms}}
 \sum_{x\in{}_KS(\F_{p^m})}\operatorname{trace}\xi_x(\Phi^m).
\]
这两个公式中，$\Phi$ 均为 Frobenius；为方便，用 $\xi_x(\Phi^m)$ 表示 $\Phi^m$ 在其不动点 $x$ 处 $E_\xi$ 的茎（stalk）上的作用。另一方面，可以定义类似第 1 节 $m(\pi_\R)$ 的数 $m(\pi_\R,\xi)$，例如它满足除非 $\operatorname{Ext}_{(\mathfrak g,\mathfrak k)}^*(\widetilde\xi,\pi_\R)\ne0$，否则 $m(\pi_\R,\xi)=0$；并令 $m(\pi)=m(\pi_\R)m(\pi^p)$。主要定理推广为
\[
 L_p(s,{}_KS,\xi)=\prod_{\pi\text{ 在 }\mathcal A_0\text{ 中}}
 L(s-n/2,\pi_p,\rho)^{m(\pi,\xi)}.
\]
证明几乎与 $\xi$ 平凡时完全相同：使用推广的迹公式，计算从 $\xi$ 经 $L^2$-诱导所得的 $\mathcal G$ 表示上的 $R_f$（$f\in C_c^\infty(\mathcal G)$）的迹：
\[
 \operatorname{trace}R_f=\sum_{\{\gamma\}}
 \operatorname{meas}(\Gamma_\gamma\backslash\mathcal G_\gamma)
 \Tr\xi(\gamma)\int_{\mathcal G_\gamma\backslash\mathcal G}f(g^{-1}\gamma g)\,dg.
\]
还使用 Langlands 的一个类似于上文所用的结果，它不但给出 ${}_KS(\overline\F_p)$ 的结构，也给出其上的 $E_\xi$ 的结构；所得结构与这一对应相容。

现在谈第二点。Langlands 在 [19] 中本质上已经表明，对应于分解 $\mathcal A_0=\bigoplus\pi$，有分解
\[
 H^*({}_KS(\C),\overline\Q)=\bigoplus H_\pi^*.
\]
只有 $m(\pi)\ne0$ 的 $\pi$ 有贡献；而对于出现的 $\pi=\pi_\R\otimes\pi_f$，表示 $\pi_f$ 可定义在 $\overline\Q$ 上。与此分解对应，$H^*({}_KS,\overline\Q_l)$ 也分解，甚至 $H^*({}_KS\times\overline\F_p,\overline\Q_l)$ 也分解。主要定理的另一个推广是，对好素数 $p$，有
\[
 L_p(s,\pi,\rho)=\prod_{i=0}^{2n}
 \det\bigl(I-(\Phi|H_\pi^i(\overline\Q))p^{-s}\bigr)^{(-1)^i}.
\]
\noindent\emph{译注：此处指数亦按所供扫描本保留；关于交错指数与 $L$、$\zeta$ 函数的约定，参见前一译注。}\footnote{\textbf{译者注（本次校订）：}除交错指数外，本式还须保留移位和重数。按定理 2.2.1 及 $H_\pi^*$ 为整个 $\pi$-等型部分的约定，应写为 $\prod_i\det(I-\Phi p^{-s}|H_\pi^i)^{(-1)^{i+1}}=L_p(s-n/2,\pi,\rho)^{m(\pi,K)}$。Frobenius 的作用应在相应的 $\overline\Q_l$ 上同调上理解。}

证明方法是：(1) 注意子空间 $H_\pi^*(\overline\Q)$ 由 $\mathcal H(G(\A_f^p),K^p)$ 中算子的作用在 $H^*({}_KS,\overline\Q)$ 内确定；(2) 应用上述论证，但允许 $f^p$ 为此 Hecke 代数的任意元素。Piatetskii-Shapiro 在 [24] 中对 $F=\Q$ 正是利用同余关系证明这个结果。

\origsubsection{6.3}{}
我也注意到，第 5 节关于 $\GL_2$ 和 $H^\times$ 上测度的结果，本质上说 Jacquet--Langlands 给出的测度，是 Serre 在 [25] 中所称的各群上的 Euler--Poincaré 测度（Euler--Poincaré measure）乘以同一个常数。严格地说，这并不完全正确，因为这些群上的 Euler--Poincaré 测度为零；但 $\SL_2$ 与 $N_H^1$ 上的此类测度之间存在明确关系。这一观察能否推广？对于实群，我预期会像 Harder 的工作 [13] 那样，出现因子 $\operatorname{card}(W_{G(\C)}/W_K)$。实际上，我怀疑当 $G$ 为各向异性群的内扭曲时，Harder 的论文正是在本质上将测度的内扭曲与 Serre 的测度联系起来。如果 $G$ 不是这样的内扭曲，Serre 的测度就恒为零；此时应如何描述测度之间的关系？

至于轨道积分的比较，Drinfeld 对 $\GL_n$ 与 $n$ 次除代数有一个很好的结果，见本研究班 Kazhdan 的报告。Diana Shelstad 已在她的学位论文 [27] 中证明了关于实群的许多漂亮结果。

\begin{thebibliography}{29}
\bibitem[1]{c29-1} T. Asai, \emph{On certain Dirichlet series associated with Hilbert modular forms and Rankin's method}, Math. Ann. 226 (1977), 81--94.
\bibitem[2]{c29-2} A. Borel, \emph{Automorphic L-functions}, 本论文集第二册，第 27--61 页。
\bibitem[3]{c29-3} A. Borel and N. Wallach, \emph{Cohomology of discrete subgroups of semisimple groups}, Notes, Inst. Advanced Study, Princeton, N.J., 1977.
\bibitem[4]{c29-4} P. Cartier, \emph{Lectures on representations of $p$-adic groups: A survey}, 本论文集第一册，第 111--155 页。
\bibitem[5]{c29-5} W. Casselman, \emph{Characters and Jacquet modules}, Math. Ann. 230 (1977), 101--105.
\bibitem[6]{c29-6} P. Deligne, \emph{Formes modulaires et représentations $l$-adiques}, Sém. Bourbaki 1968/69, no. 355, Lecture Notes in Math., vol. 179, Springer, New York.
\bibitem[7]{c29-7} P. Deligne, \emph{Travaux de Shimura}, Sém. Bourbaki 1970/71, no. 389, Lecture Notes in Math., vol. 244, Springer, New York.
\bibitem[8]{c29-8} P. Deligne, \emph{Le support du caractère d'une représentation supercuspidale}, C. R. Acad. Sci. Paris Sér. A-B 283 (1976), A155--A157.
\bibitem[9]{c29-9} V. G. Drinfeld, \emph{Proof of the global Langlands conjecture for GL(2) over a function field}, Functional Anal. Appl. 61 (3) (1977), 74--75（俄文）。
\bibitem[10]{c29-10} M. Duflo and J.-P. Labesse, \emph{Sur la formule des traces de Selberg}, Ann. Sci. École Norm. Sup. 4 (1971), 193--284.
\bibitem[11]{c29-11} M. Eichler, \emph{Introduction to the theory of algebraic numbers and functions}, Academic Press, New York, 1966.
\bibitem[12]{c29-12} S. Gelbart and H. Jacquet, \emph{Forms of GL(2) from the analytic point of view}, 本论文集第一册，第 213--251 页。
\bibitem[13]{c29-13} G. Harder, \emph{A Gauss-Bonnet formula for discrete arithmetically defined groups}, Ann. Sci. École Norm. Sup. 4 (1971), 409--455.
\bibitem[14]{c29-14} Y. Ihara, \emph{Hecke polynomials as congruence $\zeta$-functions in elliptic modular case}, Ann. of Math. (2) 85 (1967), 267--295.
\bibitem[15]{c29-15} H. Jacquet, \emph{Automorphic forms for GL(2) II}, Lecture Notes in Math., vol. 278, Springer, New York.
\bibitem[16]{c29-16} H. Jacquet and R. P. Langlands, \emph{Automorphic forms for GL(2)}, Lecture Notes in Math., vol. 114, Springer, New York, 1973.
\bibitem[17]{c29-17} R. Kottwitz, \emph{Combinatorics and Shimura varieties mod p}, 本论文集第二册，第 185--192 页。
\bibitem[18]{c29-18} S. Lang, \emph{$\SL_2(\R)$}, Addison-Wesley, 1975.
\bibitem[19]{c29-19} R. P. Langlands, \emph{Modular forms and $l$-adic representations}, Lecture Notes in Math., vol. 349, Springer, New York, 1973.
\bibitem[20]{c29-20} R. P. Langlands, \emph{Some contemporary problems with origins in the Jugendtraum}, Proc. Sympos. Pure Math., vol. 28, Amer. Math. Soc., Providence, R.I., 1976, pp. 401--418.
\bibitem[21]{c29-21} R. P. Langlands, \emph{Shimura varieties and the Selberg trace formula}, Canad. J. Math.（待刊）。
\bibitem[22]{c29-22} R. P. Langlands, \emph{On the zeta-functions of some simple Shimura varieties}（待刊）。
\bibitem[23]{c29-23} Y. Matsushima and G. Shimura, \emph{On the cohomology groups attached to certain vector-valued differential forms on the product of upper half-planes}, Ann. of Math. (2) 78 (1963), 417--449.
\bibitem[24]{c29-24} I. I. Piatetskii-Shapiro, \emph{Zeta functions of modular curves}, Lecture Notes in Math., vol. 349, Springer, New York, 1973.
\bibitem[25]{c29-25} J.-P. Serre, \emph{Cohomologie des groupes discrets}, Prospects in Mathematics, Ann. of Math. Studies, no. 70, Princeton Univ. Press, Princeton, N.J., 1971.
\bibitem[26]{c29-26} D. Shelstad, \emph{Orbital integrals for $\GL_2(\R)$}, 本论文集第一册，第 107--110 页。
\bibitem[27]{c29-27} D. Shelstad, \emph{Some character relations for real reductive algebraic groups}, Thesis, Yale Univ., 1974.
\bibitem[28]{c29-28} G. Shimura, \emph{Introduction to the arithmetic theory of automorphic functions}, Princeton, 1971.
\bibitem[29]{c29-29} G. Shimura, \emph{On the holomorphy of certain Dirichlet series}, Proc. London Math. Soc. 31 (1975), 79--98.
\end{thebibliography}
\begin{flushleft}\small University of British Columbia\end{flushleft}
